\documentclass[11pt]{article}
\usepackage[T1]{fontenc}
\usepackage{lmodern}
\usepackage[margin=1in]{geometry}
\usepackage{amsmath,amssymb,amsthm,mathtools}
\usepackage{booktabs,array,longtable}
\usepackage{enumitem}
\usepackage{xcolor}
\usepackage{microtype}
\usepackage[colorlinks=true,linkcolor=blue!55!black,citecolor=blue!55!black,urlcolor=blue!55!black]{hyperref}
\usepackage[nameinlink,noabbrev]{cleveref}

\hypersetup{
  pdftitle={A Candidate Proof of Finite Image for Rank-Three Surface-Group Representations in Genus at Least Five},
  pdfauthor={Robert Sneiderman},
  pdfsubject={Unrefereed candidate manuscript},
  pdfkeywords={mapping class groups, surface groups, local systems, representation theory}
}

\newtheorem{theorem}{Theorem}[section]
\newtheorem{proposition}[theorem]{Proposition}
\newtheorem{lemma}[theorem]{Lemma}
\newtheorem{corollary}[theorem]{Corollary}
\newtheorem{claim}[theorem]{Claim}
\theoremstyle{definition}
\newtheorem{definition}[theorem]{Definition}
\newtheorem{remark}[theorem]{Remark}
\newtheorem{audit}[theorem]{Audit point}

\newcommand{\C}{\mathbf C}
\newcommand{\Ocal}{\mathcal O}
\newcommand{\Mcal}{\mathcal M}
\newcommand{\Ecal}{\mathcal E}
\newcommand{\Lcal}{\mathcal L}
\newcommand{\Vcal}{\mathcal V}
\newcommand{\Ucal}{\mathcal U}
\newcommand{\Wcal}{\mathcal W}
\newcommand{\Hcal}{\mathcal H}
\newcommand{\Scal}{\mathcal S}
\newcommand{\Kcal}{\mathcal K}
\newcommand{\Ncal}{\mathcal N}
\newcommand{\Dcal}{\mathcal D}
\newcommand{\Bcal}{\mathcal B}
\newcommand{\Mod}{\operatorname{Mod}}
\newcommand{\rk}{\operatorname{rk}}
\newcommand{\degp}{\operatorname{pardeg}}
\newcommand{\ad}{\operatorname{ad}}
\newcommand{\im}{\operatorname{im}}
\newcommand{\Pic}{\operatorname{Pic}}
\newcommand{\Hom}{\operatorname{Hom}}
\newcommand{\End}{\operatorname{End}}
\newcommand{\tr}{\operatorname{tr}}
\newcommand{\Gr}{\operatorname{Gr}}
\newcommand{\HN}{\operatorname{HN}}
\newcommand{\SL}{\operatorname{SL}}
\newcommand{\GL}{\operatorname{GL}}
\newcommand{\PGL}{\operatorname{PGL}}
\newcommand{\SO}{\operatorname{SO}}
\newcommand{\Spec}{\operatorname{Spec}}
\newcommand{\Span}{\operatorname{Span}}
\newcommand{\coker}{\operatorname{coker}}

\title{A Candidate Proof of Finite Image\\
for Rank-Three Surface-Group Representations\\
in Genus at Least Five}
\author{Robert Sneiderman}
\date{Version 0.1.5-candidate\\19 August 2026}

\begin{document}
\maketitle

\begin{abstract}
We give a candidate proof that a complex rank-three representation of an
orientable genus-$g$ punctured surface group, whose conjugacy class has
finite mapping-class-group orbit, has finite image whenever $g\ge5$.
Landesman--Litt prove this under the strict inequality
$r<\sqrt{g+1}$, which treats rank three directly only for $g\ge9$.
The proof combines a corrected endpoint estimate, explicit
Harder--Narasimhan analysis in genera six and five, universal-extension
and boundary-cocycle obstructions for adjoint bundles, and a
rank-three-specific formal-propagation argument.  Separate
spectral-projector and canonical-Deligne determinant reconstructions are
retained for auditability but are not needed for the main elimination.  The
manuscript is unrefereed, and the assembled argument has not received
independent mathematical review.
\end{abstract}

\begin{center}
\fbox{\parbox{0.92\textwidth}{
\textbf{Status and scope.}  This is an unrefereed candidate proof.
The published inputs are recorded in \cref{sec:dependencies}.  No claim is
made in genus three or four.}}
\end{center}

\begin{center}
\small Copyright \textcopyright\ 2026 Robert Sneiderman.  Licensed under
\href{https://creativecommons.org/licenses/by/4.0/}{CC BY 4.0}.
\end{center}

\tableofcontents

\section{Statement and proof architecture}

Let $\Sigma_{g,n}$ be an orientable genus-$g$ surface with $n$ punctures.
The mapping class group $\Mod_{g,n}$ acts on conjugacy classes of
representations of $\pi_1(\Sigma_{g,n})$ by precomposition.

\begin{theorem}[Main theorem]\label{thm:main}
Let $g\ge5$ and $n\ge0$.  If
\[
 \rho:\pi_1(\Sigma_{g,n})\longrightarrow \GL_3(\C)
\]
has finite $\Mod_{g,n}$-orbit up to conjugacy, then $\rho$ has finite image.
\end{theorem}

Landesman--Litt prove finite image when
\[
 r<\sqrt{g+1},
\]
with strict inequality \cite[Theorem 1.2.1]{LLCan}.  Thus their theorem treats
$r=3$ directly for $g\ge9$.  The new cases are $g=8,7,6,5$.

The proof has eight layers.
\begin{enumerate}[label=\textup{(\arabic*)}]
\item At the square endpoint $r^2-1=g$, equality in the
Landesman--Litt bundle estimate is shown to be impossible.
\item A universal-extension construction produces large trivial subbundles
of quotient normalizers.  A Cech boundary-cocycle identity rules out such a
subbundle once it has generic rank at least three and the full extension
boundary is injective.
\item In genus five, proper positive-dimensional projective closures reduce
to adjoint constituents of ranks $5+3$ or $6+2$; the rank-six monomial
constituent is treated separately.
\item The dense rank-eight genus-five case is reduced to a finite HN branch
tree.
\item Every no-high-slope branch is eliminated by Brill--Noether, Kirillov,
relative-line, boundary-cocycle, block-vanishing, or exact jet arguments.
In particular, after closed descent the zero-weight $q=9$ branch supplies a
trivial quotient-normalizer subbundle of rank at least three.
\item The sole high-HN branch is another instance of the same obstruction:
the universal extension supplies a trivial rank-four subbundle, while
parabolic stability makes the full boundary injective.
\item For independent reconstruction, the former residual $q=9$
configuration is also eliminated by a spectral projector, and an integral
periodic-chain calculation extends the minimal-nilpotent determinant
character without poles for arbitrary contact order.
\item Projective-extension uniqueness descends residual adjoint invariants
from any dominant versal base to a finite etale cover, and a direct
small-extension devissage propagates their vanishing through every Artin
thickening when the unitary anchor is irreducible.  Reducible unitary anchors
are treated separately by the fibral-unitary isotypic decomposition and an
exact coefficient/multiplicity inequality.  The remaining rigidity,
integrality, nonabelian-Hodge, and socle-induction steps of \cite{LLCan} then
apply.
\end{enumerate}

\paragraph{Load-bearing dependency roadmap.}
The shortest proof path has four stages.
\begin{enumerate}[label=\textup{R\arabic*.},leftmargin=*]
\item The bundle and fixed-part interfaces in
\cref{prop:coparabolic-bundle,prop:closure,lem:zero-weight-closed-descent,lem:fixedpart-rankzero,lem:coefficient-unitarity}
feed the endpoint and normalizer results
\cref{prop:endpoint-pairing,cor:endpoint,prop:first-super,prop:super-classification,prop:normalizer-boundary,prop:adjoint-universal}.
\item Those results drive the finite exclusions in
\cref{sec:rank8-g76,sec:monomial5,sec:hn5,sec:zero-weight5,sec:punctured-no-high,sec:high}.
\item The resulting adjoint vanishing enters
\cref{prop:finite-index-descent,lem:artin-devissage,lem:constant-gap,lem:reducible-multiplicity,prop:formal-propagation}.
\item The final unitary, semisimple, and socle steps are
\cref{prop:unitary-rank3,prop:semisimple-rank3,prop:rank3-socle}, which imply
\cref{thm:main}.
\end{enumerate}
Here \cref{lem:fixed-part-hodge,prop:relative-hodge-line} supplies the
relative globalization used inside the finite branch eliminations.  Neither
\cref{prop:q9-reduction,prop:q9-spectral} nor \cref{sec:deligne} occurs in
this chain; those results are independent reconstructions.

\section{Published inputs and conventions}\label{sec:dependencies}

We use the following statements exactly in the forms indicated.

\subsection{Landesman--Litt}

From \cite{LLCan}:
\begin{itemize}
\item Theorem 1.2.1: finite image for an MCG-finite rank-$r$ representation
when $r<\sqrt{g+1}$, with arbitrary punctures.
\item Proposition 2.1.1: semisimple subquotients preserve MCG-finiteness.
\item Lemma 2.2.3: an irreducible MCG-finite representation admits a
\emph{projective} globalization on a scheme finite etale over $\Mcal_{g,n}$.
\item Proposition 2.3.4 and Corollary 2.3.5: a finite-determinant linear lift
exists after a dominant etale base change.  We never use this dominant lift
as though it were a finite cover.
\item Lemmas 2.1.4--2.1.5, Lemma 2.2.3, and the proof of Lemma 2.2.2:
finite-index image of a dominant versal base, geometric existence of the
projective globalization, and uniqueness of the projective intertwiners
defining it.  We do not attribute a separate uniqueness statement to
Lemma 2.2.3.
\item Lemmas 2.3.2, 2.4.1, and 2.4.2: finite determinant, extension of
unitarity/finiteness from a normal fibre subgroup, and the published general
fibral-unitary Artin decomposition.  Direct adjoint devissage treats an
irreducible unitary anchor; Lemma 2.4.2 is essential for reducible anchors.
\item Theorem 1.7.1 and Lemma 6.1.1: a nonzero sub-local system of
$R^1\pi_*U$ has rank at least $2g-2\rk U$, and a low-rank sub-local system
produces a Hodge vector with controlled period-map rank.
\item Proposition 5.2.3 and Theorem 6.2.1: the parabolic multiplication
estimate and the published strict Artin vanishing.
\item Sections 8.2--8.7: strong rigidity, integrality, unitary finiteness,
nonabelian-Hodge constancy, and characteristic-socle induction.
\end{itemize}

From \cite{LLGeo} we use Lemma 6.2.3 and Proposition 6.3.6, including the
inequalities appearing in their proofs after quotienting the high HN part.

\subsection{Small-slope Brill--Noether theory}

From Brambila-Paz--Grzegorczyk--Newstead \cite{BPGN} we use:
\begin{itemize}
\item the vector-bundle Clifford estimate;
\item the generated-subbundle conclusion in the low-slope equality range;
\item the necessary inequality
\[
 n\le d+(n-k)g
\]
for a semistable bundle of rank $n$, degree $0<d\le n$, and at least $k$
sections.
\end{itemize}
We apply these estimates separately to HN graded pieces.

\subsection{Sections of universal curves}

Earle--Kra prove that for genus at least three the Teichmuller curve with $n$
marked points has exactly the $n$ tautological holomorphic sections
\cite{EarleKra}; for $n=0$ there are none.  Chen--Salter prove that the Birman
exact sequence does not virtually split for $g\ge4$ \cite{ChenSalter}.

\subsection{Canonical Deligne and coparabolic convention}

For a unitary local system, $E_*$ denotes the parabolic bundle associated
with the Deligne canonical extension, with weights in $[0,1)$.  The
parabolic non-GGG estimate is applied to
\[
 P_*=(E_*)^\vee\otimes\omega_C(D),
\]
where $\omega_C(D)$ carries the trivial parabolic structure.  The ordinary
bundle that enters all section counts is its coparabolic degree-zero piece.

\begin{remark}[Parabolic-stability bookkeeping]
\label{rem:stability-bookkeeping}
If $R\subset E$ is a nonzero proper saturated subbundle, write
$\operatorname{wt}_*(R)$ for the sum of its induced weights.  With canonical weights in
$[0,1)$,
\[
 \operatorname{wt}_*(R)\ge0,
 \qquad
 \degp(R)=\deg R+\operatorname{wt}_*(R)<0.
 \tag{2.1}
\]
Thus an ordinary subbundle of nonnegative degree contradicts parabolic
stability.  More generally, if $P\subset E$ has quotient-weight sum $\beta$
and $E$ has total weight $S$, then
\[
 \degp(P)=\deg P+S-\beta<0.
 \tag{2.2}
\]
Finally, a nonzero section of $P$ saturates to a line in $E$ of ordinary
degree at least zero, so the same calculation gives $H^0(C,P)=0$.  We cite
this remark whenever one of these elementary consequences is used.
\end{remark}

\begin{proposition}[Underlying ordinary bundle]\label{prop:coparabolic-bundle}
With the preceding convention,
\[
 \widehat P_0=E^\vee\otimes\omega_C.
\]
Consequently
\[
 \deg(E^\vee\otimes\omega_C)=\rk(E)(2g-2)+S,
\]
where $S=-\deg E$ is the total canonical parabolic weight.
\end{proposition}

\begin{proof}
Choose $0<\varepsilon$ smaller than every positive weight.  By the
coparabolic definition, $\widehat P_0=P_\varepsilon$.  The parabolic tensor
rule and the trivial parabolic structure on $\omega_C(D)$ give
\[
 P_\varepsilon
 = (E_*^\vee)_\varepsilon\otimes\omega_C(D)
   +(E_*^\vee)_0\otimes\omega_C.
\]
The parabolic dual formula used in \cite[Section~2]{LLGeo} is
\[
 (E_*^\vee)_\alpha=(\widehat E_{-\alpha})^\vee(-D).
\]
For small $\varepsilon$, periodicity gives
$\widehat E_{-\varepsilon}=E$.  Hence the first summand is
\[
 E^\vee(-D)\otimes\omega_C(D)=E^\vee\otimes\omega_C.
\]
Locally at a marked point, the second summand is a submodule of this first
summand: the factor $(-D)$ in the dual formula shifts precisely the
positive-weight directions.  Thus their sum is $E^\vee\otimes\omega_C$.

Moreover $\degp((E_*)^\vee)=0$, while the trivial parabolic line
$\omega_C(D)$ has parabolic degree $2g-2+n$.  Therefore the input to
\cite[Proposition~6.3.6]{LLGeo} lies exactly in its boundary case, with no
hidden marked-point divisor in the ordinary bundle.
\end{proof}

This identity also fixes the section deficit used below.  If the
Landesman--Litt multiplication map has rank zero, every section of
$E^\vee\otimes\omega_C$ lies in the contraction kernel, so the deficit
$\delta$ in the proof of \cite[Lemma~6.2.3]{LLGeo} is zero.

\section{Reduction to actual adjoint constituents}\label{sec:closure}

A reducible semisimple rank-three representation has constituents of rank at
most two and is already in the published strict range for $g\ge5$.  We may
therefore treat an irreducible rank-three system $V$.

\begin{proposition}[Projective closure alternatives]\label{prop:closure}
If the projective image of $V$ is infinite, the irreducible adjoint
constituents are multiplicity-free and occur in one of the following forms:
\[
\begin{array}{c|c}
\text{projective closure}&\ad V\\ \hline
\PGL_3\text{-dense}&U_8\\
\PGL_2\simeq\SO_3&U_5\oplus U_3\\
\text{monomial, component quotient }S_3&U_6\oplus U_2\\
\text{monomial, component quotient }C_3&
 U_3^+\oplus U_3^-\oplus\chi\oplus\chi^{-1}.
\end{array}
\]
The only coefficient ranks not already in the published strict range at
genus five are therefore $8$, $6$, and $5$, and each of those occurs once.
\end{proposition}

\begin{proof}
The connected reductive projective closure acts irreducibly on a
three-dimensional space.  If its connected semisimple part is nontrivial, it
is either $\SL_3$ in the standard representation or $\SL_2$ through
$\operatorname{Sym}^2$.  These give respectively $\mathfrak{sl}_3$ and
$\operatorname{Sym}^4\oplus\operatorname{Sym}^2$, with multiplicity one.

If the connected component is a torus, irreducibility gives three distinct
weights permuted transitively by a quotient $C_3$ or $S_3$.  The connected
projective torus is the full diagonal torus: its rational character space is a
nonzero invariant subspace of the irreducible two-dimensional standard
permutation module.  Hence the six root characters are distinct.  For $S_3$
they form one orbit, giving an irreducible rank-six off-diagonal module, while
the traceless diagonal part is the irreducible rank-two standard module.  For
$C_3$ the roots split into the two oriented three-cycles, giving nonisomorphic
rank-three modules $U_3^+,U_3^-$; the diagonal part splits into the two
distinct nontrivial characters.  All constituents are therefore
multiplicity-free.  Dimension three is prime, so there is no nontrivial
tensor-product case.
\end{proof}

The table records the actual residual coefficient systems that must be
excluded on finite etale moduli covers.  Formal propagation will not assume
that the finite-cover geometric contradiction persists after the dominant
etale base changes of \cite[Lemma 2.4.2]{LLCan}.  Instead,
\cref{prop:finite-index-descent,lem:artin-devissage} below works with the
whole adjoint: projective uniqueness pulls its residual local system back
from a finite cover, and a direct small-extension induction handles the
Artin thickness.

\section{The fixed-part interface}\label{sec:fixedpart}

\begin{lemma}[Closed descent in the dense zero-weight branch]
\label{lem:zero-weight-closed-descent}
Let $B\to\Mcal_{g,n}$ be a connected smooth finite-etale scheme cover
carrying the canonical projective globalization of an irreducible unitary,
projectively MCG-finite rank-three system.  Let $\pi:X\to B$ be the compact
universal curve, let $D=\sum_i s_i(B)$, and put $X^\circ=X\setminus D$.
Suppose that the fibral
projective monodromy factors through $\pi_1(\Sigma_g)$ and is Zariski dense in
$\PGL_3$.  Let $\Ucal$ be its adjoint local system on $X^\circ$, and assume
that $\Ucal$ is unitary on the total space, as ensured in the application by
\cref{lem:coefficient-unitarity}.  Then
\[
 H^0\!\left(B,R^1\pi_*\Ucal\right)
 \xrightarrow{\ \sim\ }
 H^0\!\left(B,R^1\pi_*^\circ\Ucal\right).
\]
Moreover, every nonzero fixed vector on the right, together with the
rank-zero Hodge line supplied by the fixed-part construction, descends to a
connected finite-etale scheme cover of $\Mcal_g$.
\end{lemma}

\begin{proof}
The projective puncture monodromies are trivial, so $\Ucal$ extends uniquely
across every marked divisor.  The relative residue sequence is
\begin{equation}
 0\longrightarrow R^1\pi_*\Ucal
 \longrightarrow R^1\pi_*^\circ\Ucal
 \longrightarrow\bigoplus_i s_i^*\Ucal(-1)
 \longrightarrow R^2\pi_*\Ucal.
 \label{eq:residue-fixed-part}
\end{equation}
In fact the last term vanishes because dense adjoint monodromy has no
fibrewise invariants, but only exactness is needed below.

Write $\Gamma=\pi_1(B)$ and let $K_i$ be the Birman point-pushing subgroup
for the $i$th mark.  Let
$\operatorname{ev}_i:K_i\to\pi_1(\Sigma_g)$ be evaluation followed by filling
the other marks.  The subgroup $\Gamma\cap K_i$ has finite index in $K_i$.
Under the canonical projective globalization, point-pushing by $\alpha$ acts
on $s_i^*\Ucal$ as
\[
 \operatorname{Ad}\!\left(\bar\rho(\operatorname{ev}_i(\alpha))\right).
\]
The lift of $\alpha$ along the boundary section $s_i$ moves the
universal-curve basepoint together with the $i$th marked point.  After filling
the marks, this diagonal lift acts on the based closed surface group by
conjugation with $\operatorname{ev}_i(\alpha)$.  The intertwiner calculation
in the proof of \cite[Lemma 2.2.2]{LLCan} therefore gives the displayed
adjoint action on $s_i^*\Ucal$.  This is a direct diagonal/filling computation;
Lemma 2.2.3 is used only for the geometric existence of the globalization.
After filling the marked points, the
evaluation image of $\Gamma\cap K_i$ has finite index in
$\pi_1(\Sigma_g)$.  Its image
under $\bar\rho$ is therefore still Zariski dense: a finite-index subgroup
of a dense subgroup of the connected group $\PGL_3$ is dense.  Consequently
\[
 H^0(B,s_i^*\Ucal)
 \subseteq
 \mathfrak{sl}_3^{\,\overline{\operatorname{Ad}\bar\rho(
 \operatorname{ev}_i(\Gamma\cap K_i))}}=0.
\]
Taking global flat sections in \eqref{eq:residue-fixed-part} proves the
displayed isomorphism.

Let $\bar\Gamma$ be the image of $\Gamma$ under the forgetful map to
$\Mod_g$; it has finite index.  On closed twisted cohomology, a forgotten
point-push acts by an inner automorphism together with the corresponding
coefficient conjugation, and hence acts trivially on group cohomology.
Explicitly, if $z$ is a cocycle and the push is represented by $a$, the
transformed cocycle is $g\mapsto z(g)+(g-1)z(a)$, differing from $z$ by a
coboundary.
Thus the monodromy action on $R^1\pi_*\Ucal$ factors through
$\bar\Gamma$.  The factored projective representation is stabilized by
$\bar\Gamma$.  Its dense image has trivial centralizer in $\PGL_3$, so the
uniqueness proof of \cite[Lemma 2.2.2]{LLCan} constructs the canonical
projective globalization over the corresponding finite cover of $\Mcal_g$
even if the closed projective representation has no global $\GL_3$-lift.
After passage to a common finite cover, its pullback to the pointed family is
the original globalization: both are defined by the same unique projective
intertwiners from the proof of \cite[Lemma 2.2.2]{LLCan}.
Their adjoint local systems are therefore identified.  The fixed vector
descends to the closed cover.  Intersecting with a level subgroup makes it a
scheme cover.  The closed adjoint globalization has trivial determinant and
irreducible unitary fibral restriction.  The exact sequence of
\cite[Lemma 2.1.5]{LLCan} makes the fibre group normal, so
\cite[Lemma 2.4.1(2)]{LLCan} makes the closed adjoint unitary on the total
space.  Apply \cite[Theorem 4.1.1]{LLCan} to
$\Ucal\oplus\Ucal^\vee$, which has the natural real structure noted there,
and take the $\Ucal$-summand.  Applying the same theorem with
$D=\varnothing$ identifies its closed $F^1$ term with
$\pi_*(E\otimes\omega_\pi)$.  For the open family, trivial puncture monodromy
identifies the weight-one term with $R^1\pi_*\Ucal$, and strictness of the
inclusion of variations of mixed Hodge structure gives
\[
 F^1R^1\pi_*^\circ\Ucal\cap R^1\pi_*\Ucal
 =F^1R^1\pi_*\Ucal.
\]
Hence the rank-zero Hodge representative lies in
\[
 H^0(C,E\otimes\omega_C),
\]
with no logarithmic residue, and descends with the fixed vector.
\end{proof}

\begin{lemma}[A rank-one fixed part gives multiplication rank zero]
\label{lem:fixedpart-rankzero}
Let $U$ be a unitary local system on the total space of a punctured versal
family over a smooth quasi-projective complex base whose classifying map is
etale.  If
\[
 H^0(B,R^1\pi_*^\circ U)\ne0,
\]
then, after possibly replacing $U$ by its complex conjugate, there is a
nonzero
\[
 v\in H^0(C,E\otimes\omega_C(D))
\]
for which the Landesman--Litt multiplication map
\[
 H^0(C,E^\vee\otimes\omega_C)
 \longrightarrow H^0(C,\omega_C^2(D))
\]
has rank zero.
\end{lemma}

\begin{proof}
A nonzero flat section spans a trivial irreducible rank-one sub-local system
$L\subset R^1\pi_*^\circ U$.  By \cite[Lemma 6.1.1]{LLCan}, after possible
complex conjugation there is a nonzero $v\in L_b\cap F^1\mathcal H_b$ whose
period derivative has rank at most $\rk L/2=1/2$.  Its rank is an integer,
so it is zero.  Equivalently, in the proof of that lemma the two-Hodge-type
case cannot occur in rank one, and the remaining case has identically zero
period map.  The Hodge--de Rham degeneration used in
\cite[Theorem 4.1.1]{LLCan} identifies $F^1$ with the logarithmic one-form
term.  Because the classifying map is etale, the cotangent map $c_b^*$ in
\cite[Theorem 5.1.6]{LLCan} is an isomorphism, so the factorization in that
theorem identifies zero period derivative with the displayed zero
multiplication map.
\end{proof}

\begin{lemma}[Compatibility for the new adjoint coefficients]
\label{lem:coefficient-unitarity}
After a finite etale base change fixing the projective-closure data, every
coefficient of rank $8$, $6$, or $5$ used below is a globally defined
irreducible unitary local system on the total family, and complex conjugation
does not alter its parabolic HN problem.
\end{lemma}

\begin{proof}
The dense adjoint $U_8$ is self-dual by the trace form.  The principal
$\PGL_2$ constituents $U_5$ and $U_3$ are the self-dual representations
$\operatorname{Sym}^4$ and $\operatorname{Sym}^2$.  The $S_3$ monomial
constituent $U_6$ is self-dual under the pairing of $E_{ij}$ with $E_{ji}$.
For a unitary system, complex conjugation is the dual, so these coefficients
are self-conjugate.  The two oriented rank-three constituents in the cyclic
monomial case are exchanged by conjugation, but both already lie in the
published strict range.

The determinants of $U_8,U_5,U_3$ are trivial.  On $U_6$ the product of the
six torus weights $t_i/t_j$ is one and the component group contributes only
a finite permutation determinant.  Thus every new coefficient has finite
determinant.  Its fibral restriction is irreducible and unitary, so
\cite[Lemma 2.4.1]{LLCan} makes the total coefficient unitary.  Finite etale
base change preserves versality.  Hence \cref{lem:fixedpart-rankzero} applies
on every finite cover required later.
\end{proof}

\begin{remark}[Why a general fibre may be used]\label{rem:general-fibre}
All branch calculations below are made after passing to a nonempty connected
finite-etale scheme cover $B\to\Mcal_{g,n}$ carrying the relevant fixed part.
Since $\Mcal_{g,n}$ is irreducible, this map is surjective and $B$ is
integral.  The fixed part supplies a Hodge section that is nonzero on every
fibre and whose multiplication map is zero.  Its generic HN type, saturated
image, evaluation ranks, and the resulting branch identities persist after
shrinking $B$ to a nonempty open subset.  Because $B\to\Mcal_{g,n}$ is
dominant and etale, that open subset meets the inverse image of every
prescribed nonempty open locus in moduli.  We may therefore test a generic
branch on a nonhyperelliptic fibre, or on a pointed fibre satisfying the
stated non-Weierstrass and nonspecial-divisor conditions.

In particular, on the genus-five pointed locus used below one may require
\[
 h^0\!\left(\Ocal_C(2p)\right)
 =h^0\!\left(\Ocal_C(3p)\right)
 =h^0\!\left(\Ocal_C(2p+r)\right)
 =h^0\!\left(\Ocal_C(2p+2r)\right)=1.
\]
These conditions define a nonempty open locus; the HN calculation is not a
specialization from one branch to another.
\end{remark}

\section{The endpoint rank-one lemma}\label{sec:endpoint}

\begin{proposition}[Strict stable pairing estimate]\label{prop:endpoint-pairing}
Let $C$ have genus $g\ge2$, let $D$ be reduced, and let $E_*$ be
parabolically stable of parabolic degree zero and rank $d\ge2$.  For a nonzero
\[
 v\in H^0(C,E\otimes\omega_C(D)),
\]
let $s$ be the rank of the multiplication map of
\cite[Proposition 5.2.3]{LLCan}
\[
 H^0(C,E^\vee\otimes\omega_C)\longrightarrow
 H^0(C,\omega_C^2(D))
\]
Then
\[
 \boxed{d\ge g-s+1.}
\]
\end{proposition}

\begin{proof}
The claim is automatic if $s\ge g$, so assume $s<g$.  The published weak
estimate gives $d\ge g-s$.  If the displayed stronger estimate failed, then
$d=g-s$.

Put $F=E^\vee\otimes\omega_C$, let $K$ be the corank-one contraction
kernel, and let $A$ be the saturation in $F$ of the evaluation image of
$H^0(K)$.  Saturation stays inside the saturated kernel $K$, and every
section of either $K$ or $A$ factors through $A$.  Hence, with
$\delta=h^0(F)-h^0(A)$,
\[
 \delta=h^0(F)-h^0(K)=s.
\]
If $c=d-\rk A$, Proposition 6.3.6 of \cite{LLGeo} gives
\[
 d\ge gc-s.
\]
Since $c\ge1$ and $d+s=g$, equality forces $c=1$.

Let $N^t$ be the HN part of $A$ with slope $>2g-2$.  The proof of the same
proposition, using $H^1(N^t)=0$ to preserve the deficit $s$, gives
\[
 \rk(F/N^t)\ge g-s=d.
\]
Since $\rk F=d$, $N^t=0$.  The intermediate inequality of
\cite[Lemma 6.2.3]{LLGeo} is
\[
 \mu(F)+2\le (2g-\mu(F))d+2s.
\]
The quotient-slope lemma gives $\mu(F)\ge2g-2$.  At $\mu(F)=2g-2$ the two
sides are equal because $d+s=g$, and increasing $\mu(F)$ increases the left
side while decreasing the right.  Therefore equality is forced.  If $S$ is
the total parabolic weight, then
\[
 \mu(F)=2g-2+\frac{S}{d},
\]
and therefore $S=0$.  The bundle $E$ is now ordinary stable of degree zero,
and $F$ is stable of slope $2g-2$: all weights lie in $[0,1)$, so total
weight zero forces every weight to be zero and the strictly increasing
parabolic flag to be the trivial one-step flag.

We can now return to the complete inequality chain in the proof of
\cite[Lemma 6.2.3]{LLGeo}.  Since $N^t=0$, the bundle $A$ satisfies its
Harder--Narasimhan hypotheses, and that proof gives
\[
 \deg F+(1-g)d
 \le \frac{\deg A}{2}+\rk A+s
 \le \frac{\mu(F)\rk A}{2}+\rk A+s.
\]
Here $\mu(F)=2g-2$, $\rk A=d-1$, and $d+s=g$.  Both extreme terms are
therefore equal to $d(g-1)$:
\[
 \deg F+(1-g)d=d(g-1)
\]
and
\[
 \frac{\mu(F)(d-1)}2+(d-1)+s
 =g(d-1)+s=d(g-1).
\]
Thus equality holds throughout.  In particular,
\[
 \deg A=\mu(F)\rk A,
\]
so $\mu(A)=\mu(F)$.  But $d\ge2$ and $c=1$, hence $A$ is a nonzero proper
subbundle of the ordinary stable bundle $F$.  Stability requires
$\mu(A)<\mu(F)$, a contradiction.

This last step uses no additional Clifford theorem: the contradiction is
already contained in the equality conditions of the Landesman--Litt
inequality chain.  The condition $d\ge2$ is exactly what ensures
$\rk A=d-1>0$.
\end{proof}

\begin{remark}[The rank-one exception]
The restriction $d\ge2$ is sharp.  If $D=\varnothing$ and
$E\in\Pic^0(C)$ is nontrivial, then
$h^0(E\omega_C)=h^0(E^{-1}\omega_C)=g-1$.  Multiplication by any nonzero
$v\in H^0(E\omega_C)$ is injective on $H^0(E^{-1}\omega_C)$, so
$s=g-1$ and $d+s=g$.  Thus the published weak bound is attained in rank
one.
\end{remark}

\begin{corollary}[Strict fixed-part estimate]\label{cor:endpoint}
Let $U$ be an irreducible unitary fibral local system of rank $d$ on a
punctured versal genus-$g$ family.  Every nonzero irreducible sub-local
system $L\subset R^1\pi_*U$ has
\[
 \rk L\ge
 \begin{cases}
 2g-2,&d=1,\\
 2g-2d+2,&d\ge2.
 \end{cases}
\]
\end{corollary}

\begin{proof}
For $d=1$ use \cite[Theorem 1.7.1]{LLCan}.  For $d\ge2$, by
\cite[Lemma 6.1.1]{LLCan}, after possible complex conjugation there is a
nonzero Hodge vector whose multiplication rank $s$ is at most $\rk L/2$.
The proposition gives $s\ge g-d+1$, and the claim follows.
\end{proof}

\begin{corollary}[Artin vanishing through rank $g$]\label{cor:artin-endpoint}
Let $g\ge3$.  In the setup of \cite[Theorem 6.2.1]{LLCan}, the conclusion
remains valid when $\rk_A\Vcal\le g$.
\end{corollary}

\begin{proof}
A possible term $U\otimes\pi^*W$ satisfies
\[
 uw\le g
\]
and, if it contributes a nonzero invariant, the preceding fixed-part lower
bound is at most $w$.  If $w=1$, that lower bound is at least two for every
$u\le g$.  If $w\ge2$ and $u=1$, then
\[
 uw=w\ge2g-2>g.
\]
If $w\ge2$ and $u\ge2$, then $u\le g/2$ and
\[
 uw\ge2u(g-u+1)\ge2g>g.
\]
Both alternatives contradict $uw\le g$.
\end{proof}


\begin{theorem}[Square-endpoint finite-image theorem]\label{thm:square-endpoint}
Let $g\ge3$, $n\ge0$, and $r^2\le g+1$.  Every MCG-finite representation
\[
 \rho:\pi_1(\Sigma_{g,n})\longrightarrow\GL_r(\C)
\]
has finite image.
\end{theorem}

\begin{proof}
Only the boundary case $r^2=g+1$ is new.  Then
$\rk(\ad\rho)=r^2-1=g$, so \cref{cor:artin-endpoint} replaces
\cite[Theorem 6.2.1]{LLCan} in the proofs of strong cohomological rigidity
\cite[Proposition 8.2.1]{LLCan} and formal constancy
\cite[Lemma 8.5.1]{LLCan}.  The projective and linear integrality arguments,
unitary finiteness, and the semisimple nonabelian-Hodge deformation in
\cite[Sections 8.3--8.5]{LLCan} then apply unchanged.  The PVHS-to-unitary
input is used in the larger range $r<2\sqrt{g+1}$, which is automatic here.

It remains to check the non-semisimple step at equality.  In the notation of
\cite[Lemma 8.6.1]{LLCan}, let the characteristic semisimple factors have
ranks $a,b$.  Then
\[
 ab\le\frac{(a+b)^2}{4}\le\frac{g+1}{4}.
\]
Writing
$\rho_2^\vee\otimes\rho_1=\bigoplus_i\sigma_i^{\oplus n_i}$ gives
$n_i\dim\sigma_i\le(g+1)/4$ and $\dim\sigma_i<g$.  Any nonzero projected
extension support has, by \cite[Theorem 7.2.1]{LLCan}, rank at least
\[
 2g-2\dim\sigma_i
 \ge\frac{3g-1}{2}
 >\frac{g+1}{4}
 \ge n_i\dim\sigma_i,
\]
a contradiction.  Thus every characteristic extension splits.  The socle
induction of \cite[Section 8.7]{LLCan} completes the proof.
\end{proof}

This includes the square endpoints $(g,r)=(3,2),(8,3),(15,4),(24,5),\ldots$.
The detailed endpoint proof is logically independent of the genus-seven,
six, and five branch analysis below.

\section{The first super-endpoint}\label{sec:first-super}

The next proposition moves one rank above the stable endpoint.  It is used
only when the multiplication rank is zero, but it replaces two separate
low-genus calculations.

\begin{proposition}[First super-endpoint pairing lemma]
\label{prop:first-super}
Let $g\ge3$, let $D$ be reduced, and let $E_*$ be parabolically stable of
parabolic degree zero and rank
\[
 d=g+1.
\]
Write $S$ for the total canonical parabolic weight of $E_*$.  Suppose
$S\ne1$.  Then no nonzero
\[
 v\in H^0(C,E\otimes\omega_C(D))
\]
can have multiplication rank zero.
\end{proposition}

\begin{proof}
Assume the multiplication map is zero and put
$F=E^\vee\otimes\omega_C$.  Let $K\subset F$ be the contraction kernel and
let $A\subset K$ be the saturation of the evaluation image.  Then
$h^0(A)=h^0(F)$.  If
\[
 c=d-\rk A\ge1,
\]
the proof of \cite[Proposition 6.3.6]{LLGeo} gives $d\ge gc$.  Since
$g+1<2g$, one has $c=1$ and hence $\rk A=g$.

Let $N\subset A$ be the HN part of slope $>2g-2$.  The same proof, after
quotienting by $N$, gives $\rk(F/N)\ge g$, so $\rk N\le1$.  If $\rk N=1$,
then $\rk(F/N)=g$ and the complete inequality chain in
\cite[Lemma 6.2.3]{LLGeo} is
\[
 2g\le\mu(F/N)+2
 \le(2g-\mu(F/N))g
 \le2g.
\]
Thus equality holds throughout,
\[
 \mu(F/N)=\mu(A/N)=2g-2.
\]
Since
\[
 \deg F=(g+1)(2g-2)+S,
\]
this forces
\[
 \deg N=2g-2+S,
 \qquad
 \deg(A/N)=(g-1)(2g-2),
\]
and therefore $\deg(F/A)=2g-2$.  If $M\subset E$ is the saturated line
generated by $v$, contraction gives
\[
 F/A\simeq M^\vee\otimes\omega_C,
\]
so $\deg M=0$.  Its induced parabolic weights are nonnegative, contradicting
parabolic stability.  Hence $N=0$.

The intermediate inequality now gives
\[
 \mu(F)+2\le(2g-\mu(F))(g+1).
\]
Using
\[
 \mu(F)=2g-2+\frac{S}{g+1}
\]
yields
\[
 S\le\frac{2(g+1)}{g+2}<2.
\]
The parabolic-degree identity makes $S$ a nonnegative integer.  Since
$S\ne1$, it follows that $S=0$.  Thus $E$ is an ordinary stable
zero-degree bundle and $F$ is stable of slope $2g-2$.  Since $g+1>1$,
$H^0(E)=0$, and Riemann--Roch gives
\[
 h^0(F)=(g+1)(g-1)=g^2-1.
\]

The generically globally generated bundle $A$ has rank $g$ and every HN
slope lies in $[0,2g-2)$.  The BPGN Clifford estimate and stability give
\[
 2g^2-2g-2\le\deg A<2g^2-2g,
\]
so
\[
 \deg A\in\{2g^2-2g-2,\ 2g^2-2g-1\}.
\]
Set $G=A^\vee\otimes\omega_C$.  Every HN slope of $G$ is positive, and
Serre duality gives
\[
 (\deg G,h^0(G))=(1,g)\quad\text{or}\quad(2,g+1).
\]
The BPGN small-slope bound, applied HN-piece by HN-piece, gives at most
$g-1$ sections in either total degree.  This is a contradiction.
\end{proof}



\begin{proposition}[Classification of the weight-one boundary]
\label{prop:super-classification}
Under the hypotheses of Proposition~\ref{prop:first-super}, suppose instead
that a rank-zero vector exists.  Then necessarily $S=1$, and if $M\subset E$
is its saturated Hodge line and $A\subset E^\vee\otimes\omega_C$ its
saturated evaluation kernel, then
\[
 A\simeq\omega_C^{\oplus g},\qquad \deg M=-1,
\]
and the underlying bundle is the universal extension
\[
 0\longrightarrow M\longrightarrow E
 \longrightarrow\Ocal_C^{\oplus g}\longrightarrow0.
\]
Equivalently, the boundary map $\C^g\to H^1(C,M)$ is an isomorphism.
\end{proposition}

\begin{proof}
The proof of Proposition~\ref{prop:first-super} gives $S<2$ and excludes
$S=0$, so $S=1$.  Parabolic stability gives $H^0(E)=0$, and hence
\[
 h^0(F)=\chi(F)=(g+1)(g-1)+1=g^2.
\]
The rank-$g$ bundle $A$ carries all these sections and has every HN slope in
$[0,2g-2]$.  Clifford and the HN upper bound give
\[
 g^2\le\frac{\deg A}{2}+g,
 \qquad
 \deg A\le g(2g-2),
\]
so $\deg A=2g(g-1)$.  Thus
$G=A^\vee\otimes\omega_C$ has degree zero, nonnegative HN slopes, and
\[
 h^0(G)=h^1(A)=g.
\]
A semistable degree-zero bundle with as many sections as its rank is trivial,
so $G\simeq\Ocal_C^g$ and $A\simeq\omega_C^g$.

Since $\deg F=(g+1)(2g-2)+1$, one has $\deg(F/A)=2g-1$.  The contraction
identity $F/A\simeq M^\vee\otimes\omega_C$ gives $\deg M=-1$.  Tensoring the
contraction sequence by $\omega_C^{-1}$ and dualizing gives
\[
 0\to M\to E\to\Ocal_C^g\to0.
\]
The boundary map $\C^g\to H^1(M)$ is injective because
$H^0(E)=H^0(M)=0$, and both spaces have dimension $g$.  Hence it is an
isomorphism.
\end{proof}

\begin{proposition}[Normalizer boundary-cocycle obstruction]
\label{prop:normalizer-boundary}
Let $C$ be a smooth projective connected complex curve.  Let $E$ be a
vector bundle of Lie algebras whose generic fibre is
$\mathfrak{sl}_3$, let $M\subset E$ be a saturated line, and put
\[
 G=E/M,\qquad K=N_E(M)/M\subset G.
\]
Let
\[
 \partial:H^0(C,G)\longrightarrow H^1(C,M)
\]
be the boundary map.  If $\partial$ is injective, there is no vector
subspace $U\subset H^0(C,K)$ of dimension at least three for which
\[
 U\otimes\Ocal_C\longrightarrow K
\]
has generic rank $\dim U$.
\end{proposition}

\begin{proof}
Since $M$ is saturated, $G$ is locally free.  The quotient bracket identifies
\[
 K=\ker\bigl(G\longrightarrow G\otimes M^{-1}\bigr)
   =N_E(M)/M.
\]
In particular, $K$ is saturated in $G$.  The normalizer is a Lie
subalgebra, $M$ is an ideal in it, and hence $K$ is a sheaf of Lie algebras.

The action of the normalizer on $M$ induces the grading character
\[
 \lambda:K\longrightarrow\End(M)=\Ocal_C,
 \qquad [x,m]=\lambda(x)m.
\]
For a global section $x$ of $K$, the function $\lambda(x)$ is constant
because $C$ is projective.

Let $x,y\in H^0(C,K)$.  On an affine cover choose local lifts
$X_i,Y_i\in N_E(M)$ and write
\[
 a_{ij}=X_j-X_i\in M,\qquad b_{ij}=Y_j-Y_i\in M.
\]
With the corresponding Cech convention, these cocycles represent
$\partial x$ and $\partial y$.  Since $[M,M]=0$,
\[
 [X_j,Y_j]-[X_i,Y_i]
 =\lambda(x)b_{ij}-\lambda(y)a_{ij}.
\]
Consequently
\[
 \partial[x,y]
 =\lambda(x)\partial y-\lambda(y)\partial x.
\]
Injectivity of $\partial$ gives the identity
\begin{equation}
 [x,y]=\lambda(x)y-\lambda(y)x.
 \label{eq:normalizer-boundary-law}
\end{equation}
Thus every vector subspace of $H^0(C,K)$ is bracket-closed.

Suppose now that $U$ has generic rank at least three.  Replacing it by a
three-dimensional subspace, let $F=\mathbb C(C)$ and let
$U_F\subset K_F$ be its three-dimensional generic fibre.  The quotient
boundary law \eqref{eq:normalizer-boundary-law} extends $F$-bilinearly to
$U_F$.  The quotient
normalizer of a nonzero element of $\mathfrak{sl}_3$ has dimension
$1,2,3,$ or $4$ in the regular nonnilpotent, regular nilpotent,
semisimple $(2,1)$, or minimal-nilpotent cases, respectively.  Hence only
the last two cases remain.

In the semisimple $(2,1)$ case,
\[
 K_F\simeq\mathfrak{sl}_2(F),\qquad \lambda=0.
\]
Since $\dim_F U_F=3$, one has $U_F=K_F$, whereas
\eqref{eq:normalizer-boundary-law} makes it abelian, a contradiction.

In the minimal-nilpotent case, after extending $F$ if necessary, take
$m=E_{12}$.  Modulo $Fm$, the quotient normalizer has basis
\[
 D=\operatorname{diag}(1/2,-1/2,0),\quad
 H=\operatorname{diag}(1/2,1/2,-1),\quad
 P=E_{13},\quad Q=E_{32},
\]
with
\[
 \lambda(D)=1,\qquad \lambda(H)=\lambda(P)=\lambda(Q)=0
\]
and
\[
 [D,P]=\tfrac12P,\quad [D,Q]=\tfrac12Q,\quad
 [H,P]=\tfrac32P,\quad [H,Q]=-\tfrac32Q,\quad [P,Q]=0.
\]
If $\lambda|_{U_F}=0$, then
$U_F=\ker\lambda=\langle H,P,Q\rangle_F$, but this is not abelian.
If $\lambda|_{U_F}\ne0$, choose $X\in U_F$ with $\lambda(X)=1$ and put
$V=U_F\cap\ker\lambda$.  Equation
\eqref{eq:normalizer-boundary-law} gives
\[
 [V,V]=0,\qquad [X,v]=v\quad(v\in V).
\]
If two independent vectors
$u=\alpha H+\beta P+\gamma Q$ and
$v=\alpha' H+\beta' P+\gamma' Q$ commute, their $P$- and $Q$-coefficients
give
\[
 \alpha\beta'-\alpha'\beta=0,
 \qquad
 \alpha\gamma'-\alpha'\gamma=0.
\]
If either $\alpha$ or $\alpha'$ were nonzero, these equations would make
$u,v$ proportional.  Thus the only two-dimensional abelian subspace of
$\langle H,P,Q\rangle_F$ is $\langle P,Q\rangle_F$.  Writing
\[
 X=D+cH+aP+bQ,
\]
the eigenvalues of $\ad_X$ on $P,Q$ are
\[
 \frac{1+3c}{2},\qquad\frac{1-3c}{2},
\]
which cannot both equal one.  This is again a contradiction.
\end{proof}

\begin{proposition}[Adjoint universal-extension obstruction]
\label{prop:adjoint-universal}
Let $C$ have genus $g\ge5$.  Let $E$ be an adjoint bundle with generic
fibre $\mathfrak{sl}_3$, let $M\subset E$ be a saturated line of degree
$-1$, put $L=M^{-1}$ and $G=E/M$, and suppose that
\[
 W\otimes\Ocal_C\subset G,
 \qquad \dim W=g.
\]
Let $Q=G/(W\otimes\Ocal_C)$ and let $P\subset E$ be the inverse image of
$W\otimes\Ocal_C$.  Assume $H^0(P)=0$, $h^0(Q\otimes L)\le h$, and that
bracket with $M$ descends to
\[
 \beta:G\longrightarrow G\otimes L.
\]
Then
\[
 \Ocal_C^{g-h}\subset\ker\beta.
\]
Consequently, the configuration is impossible if $g-h\ge5$.  If in
addition $H^0(C,E)=0$, it is impossible already when $g-h\ge3$.
\end{proposition}

\begin{proof}
Riemann--Roch gives $h^1(M)=g$.  The boundary map of
\[
 0\to M\to P\to W\otimes\Ocal_C\to0
\]
is injective because $H^0(P)=0$, hence is an isomorphism
$e:W\xrightarrow{\sim}H^1(M)$.  Let
$\phi:W\to H^0(Q\otimes L)$ be the quotient component of $\beta$ and put
$W_0=\ker\phi$, so $\dim W_0\ge g-h$.

If $h^0(L)=0$, the remaining block
$W_0\otimes\Ocal_C\to W\otimes L$ vanishes.  If
$L=\Ocal_C(r)$, write that block as $sA$, where $s$ is the unique section
of $L$ and $A:W_0\to W$ is constant.  The morphism-of-extensions identity
says that the composite
\[
 W_0\xrightarrow{A}W\xrightarrow{e}H^1(M)
 \xrightarrow{s_*}H^1(\Ocal_C)
\]
is zero.  The last map is an isomorphism by
$0\to\Ocal_C(-r)\to\Ocal_C\to k(r)\to0$, so $A=0$.  Thus
$W_0\otimes\Ocal_C\subset\ker\beta$.

For a nonzero $m\in\mathfrak{sl}_3$, the quotient-adjoint kernel has
dimension $4,3,2,$ or $1$ according as $m$ is minimal nilpotent,
semisimple of type $2+1$, regular nilpotent, or regular semisimple.
Therefore a kernel of rank at least five is impossible.  If $H^0(C,E)=0$,
the boundary map
$H^0(C,G)\to H^1(C,M)$ is injective.  Proposition
\ref{prop:normalizer-boundary}, applied to
$\Ocal_C^{g-h}\subset\ker\beta$, gives the stronger conclusion when
$g-h\ge3$.
\end{proof}

\begin{corollary}[Two rank-three applications]
\label{cor:first-super-apps}
The rank-eight dense adjoint coefficient in genus seven and the rank-six
off-diagonal monomial coefficient in genus five have no rank-one fixed part.
\end{corollary}

\begin{proof}
For the dense adjoint in genus seven, Proposition~\ref{prop:first-super}
reduces a possible rank-zero vector to the weight-one universal extension of
Proposition~\ref{prop:super-classification}.  There $G=E/M\simeq\Ocal_C^7$,
so Proposition~\ref{prop:adjoint-universal} applies with $h=0$ and gives an
impossible rank-seven quotient-bracket kernel.

For the rank-six off-diagonal monomial coefficient in genus five,
$d=6=g+1$ and the exact local totals are
\[
 0,\quad\frac32,\quad2,\quad\frac52,\quad3.
\]
The global total is integral and therefore cannot equal one.  Apply
Proposition~\ref{prop:first-super}.
\end{proof}

\section{The rank-eight coefficient in genera seven and six}\label{sec:rank8-g76}

\subsection{Genus seven}

The Artin coefficient optimization leaves only ranks seven and eight with
base multiplicity one.  Rank seven is covered by
Corollary~\ref{cor:endpoint}, while rank eight is excluded by
Corollary~\ref{cor:first-super-apps}.  No separate genus-seven HN branch
calculation is needed.

\subsection{Genus six}

Here $F$ has rank eight, $2g-2=10$, and $\deg F=80+S$.  A rank-one fixed
part gives a saturated line $M\subset E$ and a rank-seven evaluation kernel
$A\subset F$ with
\[
 F/A\simeq M^\vee\otimes\omega_C,
 \qquad q:=\deg(F/A)=10-\deg M.
\]
Since
\[
 \degp(M)=\deg M+\operatorname{wt}_*(M)<0,
 \qquad \operatorname{wt}_*(M)\ge0,
\]
the bookkeeping in Remark~\ref{rem:stability-bookkeeping} gives $\deg M\le-1$ and
$q\ge11$.
Let $N\subset A$ be the HN part of slope $>10$, and write
$t=\rk N$, $x=\deg N$, and $y=x-S$.  The quotient-rank estimate gives
$t\le2$.

If $t=2$, then $\rk(F/N)=g=6$, so equality holds throughout the
Landesman--Litt chain.  Consequently
\[
 \mu(F/N)=\mu(A/N)=10,
 \qquad \deg(F/A)=10,
\]
and hence $\deg M=0$.  Then
$\degp(M)=\operatorname{wt}_*(M)\ge0$, contrary to
Remark~\ref{rem:stability-bookkeeping}.

Suppose $t=1$ and put $B=A/N$.  The two quotient-slope bounds give
\[
 10\le\mu(F/N)=\frac{80-y}{7}
 \le\frac{12\cdot7-2}{8}=\frac{41}{4},
\]
so $y\in\{9,10\}$.  Since
\[
 \rk B=6,
 \qquad \deg B=80-q-y,
\]
the HN upper degree $\deg B\le60$ and the Clifford inequality, together with
$q\ge11$ above, give
\[
 q+y\ge20,
 \qquad q\le y+2,
 \qquad q\ge11.
\]
Let $G=B^\vee\otimes\omega_C\subset E/M$ and let $P\subset E$ be its
inverse image.  The two offsets and all allowed values of $q$ are therefore
\[
\begin{array}{c|c|c|c}
y& q&\deg P&\deg(E/P)\\ \hline
9&11&-1&1-S\\
10&11,12&0&-S.
\end{array}
\]
The row $y=10$ contradicts parabolic stability because the proper subbundle
$P$ has ordinary degree zero and nonnegative induced parabolic weights.
For $y=9$, the inequality $x=S+9>10$ gives $S\ge2$.  Let $\beta$ be the
total parabolic weight of the quotient line $E/P$, and let $m$ be the number
of nonscalar punctures.  Stability of $P$ would require
\[
 -1+(S-\beta)<0,
 \qquad\text{hence}\qquad \beta>S-1.
\]
At a scalar puncture the adjoint parabolic structure is trivial, so its
quotient-line weight is zero.  At each nonscalar puncture that weight is
strictly less than one, while every nonscalar adjoint puncture contributes
total weight $2$ or $3$.  Consequently
\[
 \beta<m\le\frac S2\le S-1,
\]
a contradiction.  Thus no high-HN branch survives.

With $N=0$, the exact no-high list is
\[
 (S,q)=(0,11),(0,12),(0,13),(0,14),(2,12).
\]
For the zero-weight cases put $G=E/M$.  The relevant BPGN bounds, applied to all
positive-degree HN factors, are
\[
\begin{array}{c|c|c|c}
q&\deg G&h^0(G)&\text{maximum allowed}\\ \hline
11&1&6&6\\
12&2&7&6\\
13&3&8&7\\
14&4&9&8.
\end{array}
\]
Hence only $(S,q)=(0,11)$ remains among the zero-weight branches.  There the
rank-seven degree-one quotient $G$ has a saturated trivial evaluation
subbundle $W=\Ocal_C^6$.  Its quotient $Q$ is a line, $\deg M=-1$, and
$h^0(Q\otimes M^{-1})\le1$ on the nonhyperelliptic test fibre of
Remark~\ref{rem:general-fibre}.  The inverse image $P\subset E$ has $H^0(P)=0$
by the section-vanishing bookkeeping in Remark~\ref{rem:stability-bookkeeping}.
Proposition~\ref{prop:adjoint-universal},
with $g=6$ and $h=1$, therefore produces an impossible rank-five
quotient-bracket kernel.

For $(S,q)=(2,12)$ the same low-degree equality analysis gives
\[
 E/M\simeq\Ocal_C^7,
 \qquad \deg M^{-1}=2.
\]
The ordinary $\mathfrak{sl}_3$ fibre and the contracted
\[
 \mathfrak{gl}_2\ltimes(\C^2\oplus(\C^2)^\vee)
\]
fibre have no one-dimensional ideals, so the quotient bracket is fibrewise
nonzero.  On the nonhyperelliptic test fibre of
Remark~\ref{rem:general-fibre}, a degree-two line has at
most one section, and the matrix has a common zero.  Contradiction.

We have therefore obtained the required adjoint Artin vanishing for every
rank-three unitary system in genus at least six.

\section{Genus five: monomial coefficient}\label{sec:monomial5}

Let $U_6$ be the irreducible off-diagonal constituent of an irreducible
monomial rank-three unitary system.  Its rank is
\[
 6=g+1
\]
in genus five.  The exact local total weights are
\[
 0,\quad\frac32,\quad2,\quad\frac52,\quad3.
\]
The global total is integral because $\deg E+S=0$, and hence cannot equal
one.  Proposition~\ref{prop:first-super} excludes a rank-one fixed part.
The rank-five orthogonal constituent is covered by
Corollary~\ref{cor:endpoint}, and the smaller constituents lie in the
published strict range.  Thus every proper positive-dimensional projective
closure is eliminated in genus five.

\section{Genus-five dense HN reconstruction}\label{sec:hn5}

Let $E_*$ be the parabolically stable rank-eight adjoint Deligne bundle.  Put
\[
 F=E^\vee\otimes\omega_C,
 \qquad
 \deg F=64+S,
 \qquad
 h^0(F)=32+S.
\]
A rank-one fixed part gives a saturated line $M\subset E$ and a rank-seven
kernel $A\subset F$ with
\[
 F/A\simeq M^\vee\otimes\omega_C,
 \qquad
 q:=\deg(F/A)=8-\deg M.
\]
Here
\[
 \degp(M)=\deg M+\operatorname{wt}_*(M)<0,
 \qquad \operatorname{wt}_*(M)\ge0,
\]
so Remark~\ref{rem:stability-bookkeeping} gives $\deg M\le-1$ and $q\ge9$.

Let $N$ be the HN part of $A$ with slope $>8$, and put
$t=\rk N$, $x=\deg N$, $B=A/N$.  The proof of
\cite[Proposition 6.3.6]{LLGeo} and Clifford give
\begin{align}
64-(8-t)\frac{10(8-t)-2}{9-t}&\le x-S\le8t,
\label{eq:offset}\tag{7.1}\\
x&\ge S+10t+q-14,
\label{eq:cliff}\tag{7.2}\\
x&\ge8t+1.
\label{eq:high}\tag{7.3}
\end{align}
The integral offsets are
\[
\begin{array}{c|c}
t&x-S\\ \hline
1&5,6,7,8\\
2&15,16\\
3&24.
\end{array}
\]

For $t=3$, \eqref{eq:cliff} gives $q\le8$.  For $t=2$, let $P\subset E$
be the inverse image of $(A/N)^\vee\otimes\omega_C\subset E/M$.  At offset
$16$,
\[
 \rk P=6,\qquad \deg P=16-2\cdot8=0,
 \qquad \degp(P)=\operatorname{wt}_*(P)\ge0,
\]
contrary to Remark~\ref{rem:stability-bookkeeping}.  The offset $15$
forces
\[
 t=2,\quad q=9,\quad x=S+15,\quad S\ge2,
 \quad(A/N)^\vee\otimes\omega_C\simeq\Ocal_C^5.
\label{eq:highres}\tag{7.4}
\]
For $t=1$, put $y=x-S$, $B=A/N$, and
$G=B^\vee\otimes\omega_C\subset E/M$, and let $P\subset E$ be the inverse
image of $G$.  The HN upper degree and Clifford inequalities give
\[
 q+y\ge16,
 \qquad q\le y+4,
 \qquad q\ge9.
\]
Moreover
\[
 \deg P=y-8,
 \qquad \deg(E/P)=8-y-S.
\]
Thus the complete integer table is
\[
\begin{array}{c|c|c|c}
y&\text{allowed }q&\deg P&\text{extra consequence of }x=S+y>8\\ \hline
5&\varnothing&-3&--\\
6&10&-2&S\ge3\\
7&9,10,11&-1&S\ge2\\
8&9,10,11,12&0&--.
\end{array}
\]
For $y=8$, one has
$\degp(P)=\operatorname{wt}_*(P)\ge0$, contradicting
Remark~\ref{rem:stability-bookkeeping}.  For $y=6,7$, let $\beta$ be the total
weight of the quotient line $E/P$ and let $m$ be the number of nonscalar
punctures.  Stability of $P$ requires
\[
 \degp(P)=(y-8)+(S-\beta)<0,
 \qquad\text{so}\qquad
 \beta>S-2\quad(y=6),
 \qquad
 \beta>S-1\quad(y=7).
\]
At a scalar puncture the adjoint parabolic structure is trivial, so its
quotient-line weight is zero.  Every nonscalar quotient-line weight is
strictly less than one and every nonscalar adjoint puncture contributes total
weight $2$ or $3$, so $\beta<m\le S/2$.  This immediately contradicts the
$y=7$ inequality.  For
$y=6$ it also contradicts stability when $S\ge4$.  If $S=3$, there is
exactly one full-flag puncture, so $\beta<1$, whereas stability requires
$\beta>1$.  Hence every $t=1$ row is impossible, and
\eqref{eq:highres} is the only high-HN family.

If $N=0$, the intermediate inequality gives
\[
 S\le\frac{16}{3},
 \qquad
 \max(9,8+S)\le q\le14-S.
\]
Since an adjoint puncture contributes $0$, $2$, or $3$, the exact list is
\[
\begin{aligned}
S=0 &: q=9,10,11,12,13,14,\\
S=2 &: q=10,11,12,\\
S=3 &: q=11.
\end{aligned}
\]
For
\[
 G=A^\vee\otimes\omega_C\simeq E/M
\]
one has
\[
 \rk G=7,
 \qquad
 \deg G=q-8-S,
 \qquad
 h^0(G)=q-4.
\]
For the zero-weight $S=0$ branches, applying BPGN HN-piece by HN-piece gives
\[
\begin{array}{c|c|c|c}
q&\deg G&h^0(G)&\text{maximum allowed}\\ \hline
9&1&5&5\\
10&2&6&6\\
11&3&7&6\\
12&4&8&7\\
13&5&9&8\\
14&6&10&8.
\end{array}
\]
Thus only $q=9,10$ survive.  In the punctured branch $(S,q)=(2,12)$ the
rank-seven degree-two bundle has eight sections, whereas the same calculation,
now allowing degree-zero HN factors, gives a maximum of seven.  The no-high
residual list is
\[
 (0,9),(0,10),(2,10),(2,11),(3,11).
\]

\section{Zero-weight genus-five branches}\label{sec:zero-weight5}

\subsection{The branch \texorpdfstring{$(S,q)=(0,10)$}{(S,q)=(0,10)}}

Here $G=E/M$ has rank seven, degree two, and six sections, while
$L=M^{-1}=\det G$ has degree two.  The trace form $\operatorname{tr}(XY)$ (equivalently, a nonzero scalar
multiple of the Killing form) and the bracket define the Kirillov form
\[
 \Omega:\Lambda^2G\longrightarrow L.
\]
For $0\ne m\in\mathfrak{sl}_3$, the rank of $\ad_m$ is four or six.  The
radical of the induced form on $\mathfrak{sl}_3/\C m$ is the centralizer of
$m$ modulo $\C m$, so the form has the same rank.  Thus an isotropic subspace
in the seven-dimensional quotient has dimension at most five.

The positive integral HN degrees of $G$ sum to two.  The BPGN small-slope
bound shows that six sections are possible only in the following two cases:
\begin{enumerate}[label=\textup{(\roman*)}]
\item $G$ is semistable of rank seven and degree two;
\item there is an exact sequence $0\to J\to G\to Q\to0$ with
$(\rk,\deg,h^0)(J)=(1,1,1)$ and $(\rk,\deg,h^0)(Q)=(6,1,5)$, and every
section of $Q$ lifts.
\end{enumerate}

In case (i), coprimality makes $G$ stable.  Let $V$ be the saturation of the
evaluation image of its six sections.  Since $V$ is generically generated,
its HN quotients have nonnegative degree.  Stability and the degree-one BPGN
bound force $\rk V=6$ and $\deg V=0$: if $\deg V=1$, equality of the section
count would give a positive degree-one HN line in $V$, which would destabilize
$G$.  Hence the evaluation determinant has no zeros and
\[
 V\simeq\Ocal_C^6.
\]
Since $h^0(L)\le1$ on the nonhyperelliptic test fibre of
Remark~\ref{rem:general-fibre}, all entries of
$\Omega|_V$ vanish on a common fibre (or vanish identically), producing a
six-dimensional isotropic subspace.

In case (ii), let $V$ be the saturation of the evaluation image of
$H^0(Q)$.  Semistability of $Q$ of slope $1/6$ gives $\deg V=0$, while five
sections force $\rk V=5$ and
\[
 V\simeq\Ocal_C^5.
\]
Pulling $V$ back to $G$ gives $0\to J\to P\to\Ocal_C^5\to0$.  All five
sections of $Q$ lift, so this extension splits:
\[
 P\simeq J\oplus\Ocal_C^5.
\]
Let $T=Q/\Ocal_C^5$, a degree-one line.  Determinants give
$L=J\otimes T$.  The form entries on $P$ lie in $H^0(L)$ on
$\Lambda^2\Ocal_C^5$ and in $H^0(T)$ on $J\otimes\Ocal_C^5$.  If
$h^0(T)=0$, all cross terms vanish, while the $H^0(L)$ entries have a common
zero (or vanish identically).  If $h^0(T)=1$, the product of the unique
sections of $J$ and $T$ spans $H^0(L)$ whenever the latter is nonzero; the
zero of the section of $T$ kills every form entry.  In either case a fibre of
$P$ is six-dimensional isotropic, a contradiction.

\subsection{The branch \texorpdfstring{$(S,q)=(0,9)$}{(S,q)=(0,9)}}
\label{subsec:q9}

Here $S=0$ says that the projective puncture monodromies are trivial; by
itself it does not say that $n=0$.  Since this is the dense adjoint branch,
\cref{lem:zero-weight-closed-descent} moves the fixed vector and its Hodge
line to a finite cover of $\Mcal_g$.  We run the entire branch on that closed
family.  All constructions below use only the adjoint Lie bundle and its
associated Azumaya algebra, so no global $\GL_3$-lift is required.

Thus $E=\ad\bar V$ is stable of degree zero, $M\subset E$ is saturated with
$\deg M=-1$, and
\[
 G=E/M
\]
is stable of rank seven, degree one, and has five sections.  The saturated
evaluation image is
\[
 W=\Ocal_C^5\subset G,
\]
with rank-two quotient $R=G/W$, $\deg R=1$, and
$\det R\simeq L:=M^{-1}$.  Let
\[
 \beta:G\longrightarrow G\otimes L,
 \qquad K=\ker\beta.
\]
The inverse image of $W$ in $E$ is the universal extension of
$\Ocal_C^5$ by $M$.  If
\[
 h=h^0(R\otimes L),
\]
then stability and the small-slope bounds give $h\le2$, and the
universal-extension identity gives
\begin{equation}\label{eq:q9-large-kernel}
 \Ocal_C^{5-h}\subset K.
\end{equation}

Since $E$ is stable of degree zero, $H^0(C,E)=0$, so the boundary map
$H^0(C,G)\to H^1(C,M)$ is injective.  Because $h\le2$,
\eqref{eq:q9-large-kernel} supplies a generically injective trivial
subbundle of rank at least three in $K$, contradicting
Proposition~\ref{prop:normalizer-boundary}.  This eliminates the
zero-weight $q=9$ branch.

\subsection{Independent determinant and spectral reconstruction}

For auditability, we retain a second elimination of the former $h=2$
configuration.  It is not used in the proof above.  Suppose $h=2$, so
$\Ocal_C^3\subset K$.  The following proposition supplies the determinant
and degree reduction used in this independent reconstruction.

\begin{proposition}[Determinant reduction in the $h=2$ branch]
\label{prop:q9-reduction}
Let $Z\subset E$ be the saturated centralizer of the Hodge line.  One of the
following occurs.
\begin{enumerate}[label=\textup{(\roman*)}]
\item $L=M^{-1}$ is effective of degree one;
\item the Hodge line is minimal nilpotent on every fibre, the image of
$\ad_m$ is saturated, and the grading character
$\lambda:K\to\Ocal_C$ is surjective and nonzero on the displayed
$\Ocal_C^3\subset K$.
\end{enumerate}
\end{proposition}

\begin{proof}
At the generic point, $K$ is the normalizer of the Hodge line modulo that
line.  The inclusion $\Ocal_C^3\subset K$ therefore rules out every regular
element of $\mathfrak{sl}_3$: its quotient normalizer has dimension one in
the nonnilpotent case and dimension two in the regular-nilpotent case.  The
remaining nonzero orbit types are the semisimple type $(2,1)$ and the minimal
nilpotent type.

The invariant trace form identifies $Z^\perp$ with the dual quotient of $E$.
The $L$-valued Kirillov form
\[
 \kappa_m:\Lambda^2(E/Z)\longrightarrow L,
 \qquad \kappa_m(\bar x,\bar y)=\operatorname{tr}(m[x,y]),
\]
is generically symplectic, since both the semisimple $(2,1)$ orbit and the
minimal nilpotent orbit in $\mathfrak{sl}_3$ have dimension four.  Its
Pfaffian is therefore a generically nonzero morphism
\[
 \det(E/Z)\longrightarrow L^2.
\]
Because $\det E\simeq\Ocal_C$, there is a unique effective divisor $T$ such
that
\begin{equation}
 \det Z\simeq M^2\otimes\Ocal_C(T).
 \label{eq:q9-centralizer-det}
\end{equation}
The divisor $T$ is the degeneracy divisor of the integral orbit form; it
vanishes exactly when $\operatorname{im}(\ad_m)$ is saturated in
$Z^\perp\otimes L$.
Indeed, under the trace identification
$Z^\perp\simeq(E/Z)^\vee$, the determinant of
\[
 \overline{\ad}_m:E/Z\longrightarrow Z^\perp\otimes L
\]
is the square of the displayed Pfaffian.  Its divisor is therefore $2T$,
and it is a unit at every point exactly when the image is saturated.

Suppose first that the generic element is semisimple of type $(2,1)$.  Then
the generic line normalizer equals the centralizer.  More intrinsically,
bracket with $M$ defines a global grading morphism
\[
 \lambda:K\longrightarrow\Ocal_C,
 \qquad [x,m]=\lambda(x)m,
\]
whose kernel is $Z/M$.  At the generic point
$\operatorname{tr}(m^2)\ne0$, while trace invariance gives
$\lambda(x)\operatorname{tr}(m^2)=0$.  Thus $\lambda$ vanishes generically
and hence identically.  Consequently the equality extends across every
degeneration point and
\[
 K=Z/M,
 \qquad \det K\simeq M\otimes\Ocal_C(T).
\]
The inclusion $\Ocal_C^3\subset K$ has full generic rank.  Since $K$ is a
rank-three subsheaf of the stable bundle $G$ of rank seven and degree one,
$\deg K\le0$.  Taking determinants of the inclusion gives $\deg K\ge0$.
Thus $\deg K=0$, $\deg T=1$, and a nonzero section trivializes $\det K$.
Equation~\eqref{eq:q9-centralizer-det} then gives
$L\simeq\Ocal_C(T)$, so $L$ is effective.

Now suppose the generic element is minimal nilpotent.  Minimal nilpotence holds
on every fibre: the equations $m^2=0$ and $\rk m\le1$ define a closed locus,
and the saturated line never vanishes.  Let
\[
 \lambda:K\longrightarrow\Ocal_C,
 \qquad [x,m]=\lambda(x)m,
\]
be the grading character and let $Z_\lambda$ be the effective divisor defined
by $\operatorname{im}\lambda=\Ocal_C(-Z_\lambda)$.  Its kernel is
$Z/M$, and \eqref{eq:q9-centralizer-det} gives
\begin{equation}
 \det K\simeq M\otimes\Ocal_C(T-Z_\lambda).
 \label{eq:q9-kernel-det}
\end{equation}

If $\lambda$ vanishes on $\Ocal_C^3$, then
$\Ocal_C^3\subset Z/M$ has full generic rank.  Stability again gives
$\deg(Z/M)\le0$, while the determinant of the inclusion gives
$\deg(Z/M)\ge0$.  Hence $\det(Z/M)\simeq\Ocal_C$, and
\eqref{eq:q9-centralizer-det} yields $\deg T=1$ and
$L\simeq\Ocal_C(T)$.

If $\lambda$ is nonzero on $\Ocal_C^3$, its restriction is a nonzero map
from a trivial bundle to $\Ocal_C(-Z_\lambda)$.  Since $Z_\lambda$ is
effective and $C$ is projective, this forces $Z_\lambda=0$ and makes the
restriction surjective.  The integral determinant theorem of
\cref{sec:deligne} gives a nonzero morphism $\det K\to M$.  Combining it
with \eqref{eq:q9-kernel-det} produces a nonzero morphism
$\Ocal_C(T)\to\Ocal_C$, and hence $T=0$.  Thus the adjoint image is
saturated, and alternative (ii) holds.
\end{proof}

In alternative (i), \cref{prop:relative-hodge-line} gives a section of the
universal curve over the finite cover of $\Mcal_g$ supplied above,
contradicting Chen--Salter.  It remains only to eliminate alternative (ii).

\begin{proposition}[Spectral-projector obstruction]\label{prop:q9-spectral}
Suppose the Hodge line is minimal nilpotent on every fibre, $K$ has rank
four, and the grading character
\[
 \lambda:K\longrightarrow\Ocal_C,
 \qquad [x,m]=\lambda(x)m,
\]
is surjective and nonzero on the displayed $\Ocal_C^3\subset K$.  Then the
configuration is impossible.
\end{proposition}

\begin{proof}
Put $C_0=\ker\lambda$.  Then $C_0$ is the quotient centralizer
$\mathfrak z(m)/M$, of rank three, and
\[
 0\longrightarrow C_0\longrightarrow K\longrightarrow\Ocal_C\longrightarrow0.
\]
The determinant morphism gives $\deg C_0=\deg K\le-1$.  Etale-locally a
generator of the saturated minimal-nilpotent line is $m=E_{12}$.  Modulo
$M$, the centralizer then has the basis
\[
 E_{13},\quad E_{32},\quad
 H=\operatorname{diag}(1/2,1/2,-1).
\]
Its derived algebra is the span of $E_{13},E_{32}$, a saturated rank-two
subbundle.  The trace form has precisely this radical and
$\operatorname{tr}(H^2)=3/2$, so it descends perfectly to the quotient line.
Thus the derived algebra
\[
 R_0=[C_0,C_0]
\]
has rank two, and for $T=C_0/R_0$ one has
$T^{\otimes2}\simeq\Ocal_C$ and $\deg T=0$.

Let
\[
 U=\Ocal_C^3\cap C_0\simeq\Ocal_C^2.
\]
If the projection $U\to T$ vanishes, then $U\subset R_0$ and determinants
give a nonzero map $\Ocal_C\to\det R_0$, contradicting
$\deg\det R_0=\deg C_0\le-1$.  Otherwise $T\simeq\Ocal_C$.  Choose a global
section $u\in H^0(C,U)$ mapping to a generator of $T$ and a global section
$x\in H^0(C,\Ocal_C^3)$ with $\lambda(x)=1$.  Set $x_t=x+tu$.

Choose locally a generator $m$ of $M$ and a lift $X_t$ of $x_t$ to $E$.
Then $[X_t,m]=m$.  Any other lift is $X_t+fm$, and
\[
 X_t+fm=\exp(-fm)X_t\exp(fm)
\]
because $m^2=0$.  Hence the characteristic polynomial of $X_t$ is
independent of the lift and has global holomorphic, therefore constant,
coefficients.

In a frame with $m=E_{12}$, every lift has the form
\[
 \operatorname{diag}(1/2,-1/2,0)+aE_{12}+bE_{13}+cE_{32}
 +d\operatorname{diag}(1/2,1/2,-1).
\]
The section $u$ changes $d$ by a nonzero affine function of $t$.  Fix one
point of $C$ and choose $t$ so that its value $D$ is zero there.  The
three eigenvalues are
\[
 \frac{D+1}{2},\qquad \frac{D-1}{2},\qquad -D,
\]
At the chosen point the spectrum is therefore
$\{1/2,-1/2,0\}$.  Since the characteristic coefficients are global
holomorphic functions, this is the spectrum at every point; in particular,
$1/2,-1/2$ are the unique ordered pair differing by one.

Let $\alpha=1/2$, $\beta=-1/2$, and $\gamma=0$.  The projector onto the
$\gamma$-eigenspace is
\[
 P_\gamma=
 \frac{(X_t-\alpha I)(X_t-\beta I)}
      {(\gamma-\alpha)(\gamma-\beta)}.
\]
It satisfies $P_\gamma m=mP_\gamma=0$.  The conjugacy formula above makes it
independent of the choice of lift, while ordinary frame changes conjugate
it as an endomorphism.  The local projectors therefore glue to a global
rank-one idempotent in the Azumaya algebra associated with $\bar V$.  Its
trace-free part is a nonzero section of the adjoint bundle $E$, contradicting
$H^0(C,E)=0$.
\end{proof}

This completes the independent $h=2$ reconstruction.

\subsection{Fixed-part Hodge section and relative degree-one contradiction}
\label{subsec:picard}

\begin{lemma}[Fixed-part rank-zero Hodge section]
\label{lem:fixed-part-hodge}
Let $\pi^\circ:\mathcal C^\circ\to B$ be a punctured versal family over a
smooth quasi-projective complex base whose classifying map is etale, and let
$\Ucal$ be unitary.  If
\[
 H^0(B,R^1\pi_*^\circ\Ucal)\ne0,
\]
then, after complex conjugation if necessary, there is a global fibrewise
nonzero section
\[
 v\in H^0\!\left(B,\pi_*(\Ecal\otimes\omega_\pi(D))\right)
\]
whose Landesman--Litt multiplication map has rank zero on every fibre.
Equivalently, adjunction gives a global map
\[
 \omega_\pi(D)^{-1}\longrightarrow\Ecal.
\]
\end{lemma}

\begin{proof}
A nonzero fixed vector spans a trivial irreducible rank-one local system
$\mathbb L\subset R^1\pi_*^\circ\Ucal$.  For a general complex unitary
coefficient, apply \cite[Proposition 4.2.2]{LLCan} to
$\mathbb W=R^1\pi_*^\circ\widetilde{\Ucal}$ and
$\widetilde{\mathbb L}=\mathbb L$; the trivial line has its natural real
structure.  The proof of
\cite[Lemma 6.1.1]{LLCan} then splits into two cases.  Its second case would
give the real rank-one system $\widetilde{\mathbb L}=\mathbb L$ two conjugate
Hodge types, which is impossible in rank one.  Thus, after conjugating
$\Ucal$ if necessary, the first case gives a global morphism of variations
from a constant line into $F^1R^1\pi_*^\circ\Ucal$.  Its period map is
identically zero.

The filtered de Rham description and Hodge--de Rham degeneration in
\cite[Theorem 4.1.1]{LLCan} give
\[
 F^1R^1\pi_*^\circ\Ucal\simeq
 \pi_*(\Ecal\otimes\omega_\pi(D)),
\]
Because the classifying map is etale, the cotangent map $c_b^*$ in
\cite[Theorem 5.1.6]{LLCan} is an isomorphism, so the factorization in that
theorem identifies zero second fundamental form with zero multiplication
map.  Thus the corresponding global section has
rank-zero multiplication on every fibre.  The nonzero map from the
irreducible local system $\mathbb L$ is injective, so its Hodge-bundle image is
fibrewise nonzero.
\end{proof}

\begin{proposition}[Relative Hodge-line closure]
\label{prop:relative-hodge-line}
Let $\pi:X\to B$ be the pullback of the universal curve to a connected
finite-etale scheme cover, with $B$ integral and regular.  Let $\Ecal$ be a
vector bundle on $X$, and suppose a line bundle $\Hcal$ on $B$ gives a
generically nonzero morphism
\[
 \pi^*\Hcal\otimes\omega_\pi(D)^{-1}\longrightarrow\Ecal.
\]
Let $\Scal$ be the saturation of its rank-one image in $\Ecal$ and put
$\Mcal=\Scal^{**}$.  Then:
\begin{enumerate}[label=\textup{(\roman*)}]
\item $\Mcal$ is invertible and the induced map
$\Mcal\hookrightarrow\Ecal$ is injective;
\item on the generic fibre, $\Mcal_\eta$ is exactly the saturated line
generated by the Hodge vector;
\item if $\deg(\Mcal_\eta)=-1$, then $\Lcal=\Mcal^{-1}$ has relative degree
one on all of $B$;
\item if $\Lcal_\eta$ is effective, then $\pi$ admits a section.
\end{enumerate}
\end{proposition}

\begin{proof}
The total space $X$ is regular.  The sheaf $\Mcal$ is coherent reflexive of
rank one.  Since every regular local ring is a UFD, a rank-one reflexive
module on $X$ is invertible.  Double dualizing the inclusion
$\Scal\hookrightarrow\Ecal$ gives $\Mcal\to\Ecal$; it is generically
injective, and its kernel would be a torsion subsheaf of the line bundle
$\Mcal$, hence is zero.

The quotient $\Mcal/\Scal$ is supported in codimension at least two.  It does
not meet the generic fibre: a closed point of the generic curve has
codimension one in $X$.  Therefore
\[
 \Mcal_\eta=\Scal_\eta,
\]
and the right side is saturated by construction.  Fibre degree of a line
bundle is locally constant in a connected smooth proper family, proving
(iii).  The base twist $\Hcal$ has fibre degree zero.

Assume now that $\Lcal_\eta$ is effective.  Upper semicontinuity makes
\[
 \{b\in B:h^0(X_b,\Lcal_b)\ge1\}
\]
a closed subset containing the generic point, hence all of $B$.  A degree-one
line bundle on a positive-genus smooth curve has at most one independent
section, so $h^0(X_b,\Lcal_b)=1$ for every $b$; Riemann--Roch also makes
$h^1(X_b,\Lcal_b)=g-1$ constant.  Cohomology and base change therefore makes
$\pi_*\Lcal$ invertible and identifies its fibres with
$H^0(X_b,\Lcal_b)$.

The evaluation
\[
 \pi^*\pi_*\Lcal\longrightarrow\Lcal
\]
is nonzero on every fibre.  It cuts out an effective Cartier divisor
$Z\subset X$.  Its restriction to each fibre is the unique degree-one zero
divisor.  Since the section remains a nonzerodivisor on every smooth integral
fibre, $Z$ is a relative effective Cartier divisor and is flat over $B$.
It is proper with zero-dimensional fibres, hence finite.  Thus $Z\to B$ is
finite locally free of degree one, therefore an isomorphism.  Its inverse is a
section of $\pi$.
\end{proof}

Apply the proposition on the connected finite-etale base on which the fixed
part has been globalized.  Lemma~\ref{lem:fixed-part-hodge} supplies the global
evaluation, and its generic saturation is the line $M$ from the fibre
calculation.  Thus a fibrewise proof that $M^{-1}$ is effective produces an
actual section of that universal curve.  In the zero-weight effective
subcases of \cref{subsec:q9},
\cref{lem:zero-weight-closed-descent} first replaces the pointed base by a
finite cover of $\Mcal_g$; only there does the resulting section contradict
Chen--Salter's non-virtual-splitting theorem for $g\ge4$
\cite{ChenSalter}.  No line defined only on an HN stratum is being extended,
and no arbitrary pointed section is declared contradictory.

For a nontrivial rank-one sub-local system rather than an actual fixed vector,
\cite[Proposition 4.2.2 and Lemma 6.1.1]{LLCan} supply the same evaluation
with a flat line bundle pulled back from the base.  In pointed branches a
section is not itself contradictory; Earle--Kra identifies it with one of
the tautological marked sections \cite{EarleKra}.

\section{Punctured no-high branches}\label{sec:punctured-no-high}

\subsection{The branch \texorpdfstring{$(S,q)=(2,11)$}{(S,q)=(2,11)}}

Let $p$ be the unique $2+1$ puncture.  Here $G$ has rank seven,
degree one, seven sections, and nonnegative HN slopes.  The degree-one BPGN
bound shows that equality is possible only with a positive degree-one HN line
$J$ and a rank-six degree-zero quotient carrying six sections.  The quotient
is $\Ocal_C^6$, every quotient section lifts, and the extension splits.  Thus
\[
 G\simeq J\oplus\Ocal_C^6,
 \qquad
 \deg J=1,
 \qquad h^0(J)=1,
 \qquad \deg L=3.
\]
The degree-one line globalizes on the full finite cover: intrinsically
$J=\det G$, because the complementary evaluation summand has trivial
determinant.  Generic effectivity and the degree-one Abel closed immersion
give a global section, and Earle--Kra forces
$J\simeq\Ocal_C(r)$ for a marked point $r$.  Since
$\det E\simeq\Ocal_C(-2p)$,
\[
 L\simeq\Ocal_C(2p+r).
\]
The four blocks of the quotient bracket have coefficients in
\[
 \Ocal(2p+r),\quad \Ocal(2p),\quad
 \Ocal(2p+2r),\quad \Ocal(2p+r).
\]
Choose a test fibre in the pointed nonspecial-divisor locus of
Remark~\ref{rem:general-fibre}.  Each indicated section space is then
one-dimensional and every unique section vanishes at $p$.  Hence the bracket
vanishes at $p$, making $M_p$ a one-dimensional ideal of
\[
 \mathfrak{gl}_2\ltimes(\C^2\oplus(\C^2)^\vee),
\]
which has no such ideal.

\subsection{The branch \texorpdfstring{$(S,q)=(3,11)$}{(S,q)=(3,11)}}

Here $G\simeq\Ocal_C^7$ and there is one full-flag puncture $p$.
The canonical adjoint lattice has determinant $\Ocal_C(-3p)$, while
$\det G=\det E\otimes L$ is trivial.  Hence
\[
 L\simeq\Ocal_C(3p).
\]
Choose the test fibre with $p$ non-Weierstrass as permitted by
Remark~\ref{rem:general-fibre}; then
$h^0(\Ocal_C(3p))=1$, and every bracket coefficient vanishes to order at
least three at $p$.

The full-flag contracted fibre has exactly three line ideals:
\[
 \C E_{12},\qquad \C E_{23},\qquad \C E_{31}.
\]
For each, the bracket with the opposite root has exact order one:
\[
 [E_{12},E_{21}]=zH_1,
 \quad [E_{23},E_{32}]=zH_2,
 \quad [E_{31},E_{13}]=-z(H_1+H_2).
\]
These pure-root equalities also control an arbitrary saturated lift.  Let
$e_\alpha$ be the canonical root frame of one of the three ideals and
$e_{-\alpha}$ its opposite frame.  A primitive generator has the form
$m=e_\alpha+zr+O(z^2)$; multiplying by a unit removes the
$e_\alpha$ component of $r$.  Since
$[e_\alpha,e_{-\alpha}]=zH_\alpha$, while a bracket with
$e_{-\alpha}$ has Cartan component only from the $e_\alpha$ component, one
gets
\[
 \operatorname{pr}_{\mathfrak t}[m,e_{-\alpha}]
 =zH_\alpha\pmod {z^2\mathfrak t_R}.
\]
Thus the order-one Cartan term cannot cancel for any saturated lift.
A matrix of sections of $\Ocal(3p)$ would vanish to order at least three.
Contradiction.

\subsection{The branch \texorpdfstring{$(S,q)=(2,10)$}{(S,q)=(2,10)}}

Here
\[
 0\to\Ocal_C^6\to G\to D\to0,
 \qquad \deg D=0,
 \qquad \deg L=2,
\]
and $D=L(-2p)$.  If $h^0(L)=0$, then $h^0(G\otimes L)\le1$, whereas the
ordinary quotient-adjoint map has rank at least three and restricts to rank
at least two on a hyperplane.  Thus $L$ is effective.  On the
nonhyperelliptic test fibre of Remark~\ref{rem:general-fibre}, one has
$h^0(L)=1$.  Let $s$ span $H^0(L)$ and write $Z=(s)$, so
$\operatorname{length}Z=2$.

Put $W=\C^6$, and let $P\subset E$ be the inverse image of
$W\otimes\Ocal_C\subset G$.  Parabolic stability gives $H^0(E)=0$, hence
$H^0(P)=0$.  Since $M=L^{-1}$ has degree $-2$, Riemann--Roch gives
$h^1(M)=6$.  The boundary of
\[
 0\longrightarrow M\longrightarrow P
 \longrightarrow W\otimes\Ocal_C\longrightarrow0
\]
is therefore an isomorphism
\[
 e:W\xrightarrow{\sim}H^1(M).
\]

Let $Q=G/(W\otimes\Ocal_C)=D$.  The line $Q\otimes L=L^2(-2p)$ has
degree two, so nonhyperellipticity gives $h^0(Q\otimes L)\le1$.  For the
quotient component
\[
 \phi:W\longrightarrow H^0(Q\otimes L)
\]
of bracket with $M$, put $W_0=\ker\phi$.  Then $\dim W_0\ge5$.
On $W_0\otimes\Ocal_C$ the bracket lands in
$W\otimes L$.  Since $H^0(L)=\C s$, it has the form
\[
 \beta|_{W_0}=sA
 \qquad\text{for a constant map}\qquad A:W_0\longrightarrow W.
\]

The compatibility assertion is a literal morphism of extensions.  Let
$P_0\subset P$ be the inverse image of $W_0\otimes\Ocal_C$.  Bracket with
$M$, twisted by $L=M^{-1}$, gives the commutative diagram
\[
\begin{array}{ccccccccc}
0&\longrightarrow&M&\longrightarrow&P_0&\longrightarrow
 &W_0\otimes\Ocal_C&\longrightarrow&0\\
&&\downarrow 0&&\downarrow\widetilde\beta&&\downarrow sA\\
0&\longrightarrow&\Ocal_C&\longrightarrow&P\otimes L&\longrightarrow
 &W\otimes L&\longrightarrow&0.
\end{array}
\]
Indeed, $\phi|_{W_0}=0$ puts the middle image in $P\otimes L$, while
$[M,M]=0$ makes the left vertical map zero.  Pulling back the lower
extension by $sA$ and pushing out the upper extension by zero shows that its
class vanishes.  In terms of the boundary map this is exactly
\[
 s_*\circ e\circ A=0:
 W_0\longrightarrow H^1(\Ocal_C).
\]

The divisor sequence
\[
0\longrightarrow M\xrightarrow{s}\Ocal_C
\longrightarrow\Ocal_Z\longrightarrow0
\]
shows that
\[
\dim\ker\bigl(s_*:H^1(M)\longrightarrow H^1(\Ocal_C)\bigr)=1:
\]
the kernel is the cokernel of the constants
$H^0(\Ocal_C)\to H^0(\Ocal_Z)$, whose dimensions are one and two.
Since $e$ is an isomorphism, $\operatorname{rank}A\le1$, and hence
\[
 \dim\ker A\ge4.
\]
The corresponding constant sections are generically independent in
$W\otimes\Ocal_C\subset G$ and lie in
\[
 \ker\beta=N_E(M)/M.
\]
Thus $N_E(M)/M$ contains a generically rank-four trivial subbundle.  The
full boundary $H^0(G)\to H^1(M)$ is injective because $H^0(E)=0$, so
Proposition~\ref{prop:normalizer-boundary} gives a contradiction.  This
eliminates $(S,q)=(2,10)$ without identifying the divisor of $L$ with $2p$
and without a local jet calculation.

\section{The final high-HN branch}\label{sec:high}

The sole remaining branch is
\[
 t=2,\quad q=9,\quad\deg N=S+15,\quad S\ge2,
 \quad(A/N)^\vee\otimes\omega_C\simeq\Ocal_C^5.
\]
Put
\[
 G=E/M,
 \qquad L=M^{-1},
 \qquad W=\Ocal_C^5\subset G,
 \qquad Q=G/W.
\]
Then
\[
 \deg M=-1,
 \quad \deg L=1,
 \quad \rk Q=2,
 \quad \deg Q=1-S,
 \quad h^0(Q\otimes L)\le1.
\]

\subsection{Exclusion of full-flag punctures}

Let $P\subset E$ be the inverse image of $W$.  Then $\rk P=6$ and
$\deg P=-1$.  At a full-flag puncture, write the cyclic gaps as
$x,y,z>0$ with $x+y+z=1$.  The adjoint weights are two zeros and
\[
 x,y,z,1-x,1-y,1-z.
\]
Any rank-six induced fibre has weight at least the sum of the four smallest
root weights.  The six root weights total three, while each of the two
largest is strictly less than one, so the four-smallest sum is $>1$.
Therefore $\degp(P)>0$, contradicting parabolic stability.  Every non-scalar
puncture in this branch is of type $2+1$.

\subsection{Universal-extension obstruction}

Let $P\subset E$ be the inverse image of $W=\Ocal_C^5$.  The
section-vanishing bookkeeping in Remark~\ref{rem:stability-bookkeeping} gives $H^0(P)=0$,
and the branch data give
$h^0(Q\otimes L)\le1$.  Proposition~\ref{prop:adjoint-universal}, with
$g=5$ and $h=1$, yields
\[
 \Ocal_C^4\subset\ker\beta.
\]
Parabolic stability also gives $H^0(C,E)=0$: a nonzero holomorphic section
would generate a saturated parabolic line of nonnegative parabolic degree
inside the stable parabolic degree-zero bundle $E$.  Thus the full boundary
map $H^0(C,G)\to H^1(C,M)$ is injective, and
Proposition~\ref{prop:normalizer-boundary} contradicts the displayed
rank-four trivial subbundle.  Hence the final high-HN branch is impossible.

\section{Independent canonical-Deligne determinant reconstruction}
\label{sec:deligne}

The preceding boundary-cocycle obstruction makes the determinant character
unnecessary for the main proof.  We retain its integral construction and its
second elimination of the former minimal-nilpotent endpoint as an independent
reconstruction.

At a $2+1$ puncture, a scalar twist lets us choose canonical residue
\[
 \operatorname{diag}(0,0,\delta),\qquad 0<\delta<1.
\]
The adjoint residues on $E_{13},E_{23}$ are $-\delta$, those on
$E_{31},E_{32}$ are $\delta$, and the Levi part has residue zero.  Moving the
negative residues into $[0,1)$ multiplies precisely $E_{13},E_{23}$ by $z$.
Thus, with $R=\C[[z]]$ and $K=\C((z))$, the canonical adjoint lattice is
\[
 \mathfrak p_R=
 \left\{
 \begin{pmatrix}A&zb\\c&d\end{pmatrix}:\tr A+d=0
 \right\}\subset\mathfrak{sl}_3(K).
\]
Equivalently, it is the trace-zero endomorphism lattice of the periodic chain
\[
 z\Lambda_0\subset\Lambda_1\subset\Lambda_0,
 \qquad
 \Lambda_0=R^3,
 \quad
 \Lambda_1=zRe_1\oplus zRe_2\oplus Re_3.
\]
At a scalar puncture the chain is hyperspecial.  The following theorem is
stated for an arbitrary periodic chain because no contact-order or orbit
reduction is needed.

\begin{theorem}[Integral minimal-nilpotent determinant morphism]
\label{thm:integral-det}
Let $\Lambda_\bullet$ be a periodic lattice chain in an $r$-dimensional
$K$-vector space and put
\[
 \mathfrak p_R=\End_{\mathrm{ch}}(\Lambda_\bullet)\cap\mathfrak{sl}_r(K).
\]
Let $\Mcal\subset\mathfrak p_R$ be a saturated rank-one module whose generic
fibre is a minimal-nilpotent line.  If
\[
 \Kcal=
 \bigl(N_{\mathfrak{sl}_r(K)}(K\Mcal)\cap\mathfrak p_R\bigr)/\Mcal,
\]
then there is a holomorphic, generically nonzero morphism
\[
 \det\Kcal\longrightarrow\Mcal^{\otimes(r-2)}.
\]
It may have zeros at the boundary but cannot have a pole.
\end{theorem}

\begin{proof}
Let
\[
 I_K=\im m,
 \qquad H_K=\ker m,
 \qquad J_K=H_K/I_K,
 \qquad Q_K=V_K/H_K.
\]
The induced intersections and quotients of the periodic chain give free
chains $I_\bullet,J_\bullet,Q_\bullet$.  A field map in the generic
minimal-nilpotent line preserves the original chain exactly when the induced
map $Q_K\to I_K$ preserves the quotient chains.  Saturation therefore gives
the equality of integral lattices
\begin{equation}\label{eq:M-chain-final}
 \Mcal=\Hom_{\mathrm{ch}}(Q_\bullet,I_\bullet).
\end{equation}

Let $\Ncal=N_{\mathfrak{sl}_r(K)}(K\Mcal)\cap\mathfrak p_R$.  The field
normalizer preserves the flag $I_K\subset H_K\subset V_K$, hence $\Ncal$
acts on the three induced chains.  Let $\Dcal$ be its image on the graded
chains and $\Ucal$ its kernel.  There is an exact sequence
\[
 0\longrightarrow\Ucal/\Mcal\longrightarrow\Kcal
 \longrightarrow\Dcal\longrightarrow0.
\]
The two nontrivial first off-diagonal blocks give an integral injection,
generically an isomorphism,
\[
 \Ucal/\Mcal\longrightarrow
 A\oplus B,
 \qquad
 A=\Hom_{\mathrm{ch}}(J_\bullet,I_\bullet),
 \quad
 B=\Hom_{\mathrm{ch}}(Q_\bullet,J_\bullet).
\]
If both blocks vanish, the endomorphism factors from $Q_\bullet$ to
$I_\bullet$ and belongs to $\Mcal$ by \eqref{eq:M-chain-final}.  Thus
\[
 \det(\Ucal/\Mcal)\longrightarrow\det A\otimes\det B
\]
is integral and generically nonzero.

It remains to treat the diagonal factor without choosing a generic volume.
Put
\[
 \Bcal=\End_{\mathrm{ch}}(I_\bullet)\oplus
       \End_{\mathrm{ch}}(J_\bullet)\oplus
       \End_{\mathrm{ch}}(Q_\bullet).
\]
The trace map
\[
 \tau:\Bcal\longrightarrow R,
 \qquad (a,A,c)\longmapsto a+\tr(A)+c,
\]
is split surjective; its kernel is the integral trace-zero graded module
$\Dcal_0$.  Hence
\[
 \det\Dcal_0\simeq\det\Bcal
 \simeq\det\End_{\mathrm{ch}}(J_\bullet).
\]
Choose one lattice $J_0$ in the chain.  Inclusion into
$\End_R(J_0)$ and the canonical identity
$\det\End_R(J_0)\simeq R$ give an integral map
\[
 \det\Dcal_0\longrightarrow R.
\]
Since $\Dcal$ is a full-rank sublattice of $\Dcal_0$, we obtain
$\det\Dcal\to R$, again with zeros but no poles.

Finally, composition of chain maps is integral:
\[
 A\otimes B\longrightarrow\Mcal.
\]
Since $\rk J_K=r-2$, the determinant of this perfect generic pairing gives
\[
 \det A\otimes\det B\longrightarrow\Mcal^{\otimes(r-2)}.
\]
Combining the three integral determinant maps proves the theorem.
\end{proof}

\begin{corollary}[Global no-pole determinant character]
\label{cor:global-integral-det}
Let $E$ be a canonical adjoint Deligne bundle on a smooth complete curve,
with its periodic-chain lattices at the punctures.  Let $M\subset E$ be a
saturated line with generically minimal-nilpotent fibre, and let $K$ be its
global quotient normalizer.  Then the intrinsic generic determinant
character extends to a generically nonzero morphism
\[
 \det K\longrightarrow M^{\otimes(r-2)}
\]
on the whole curve.
\end{corollary}

\begin{proof}
Over the function field, the flag
$I=\im M\subset H=\ker M\subset V$ and the composition pairing
\[
 \Hom(H/I,I)\otimes\Hom(V/H,H/I)\longrightarrow\Hom(V/H,I)=M
\]
define one intrinsic determinant character.  The determinant
trivialization of the trace-zero diagonal factor is likewise functorial;
equivalently, it is the canonical identity
$\det\End(H/I)\simeq\Ocal$.  Thus every completed-DVR construction in
\cref{thm:integral-det} restricts to this same generic morphism, rather than
to unrelated local choices.  The theorem says that this rational morphism
has no pole at any closed point.  A rational section of a line bundle on a
smooth curve that is regular at every completed DVR is global, proving the
claim.
\end{proof}

\begin{proposition}[Local--global kernel identity]
\label{prop:kernel-compatibility}
Let $E$ be the canonical adjoint Deligne bundle on a smooth curve, let
$M\subset E$ be saturated, and put $G=E/M$.  If
\[
 \beta:G\longrightarrow G\otimes M^{-1}
\]
is the quotient bracket and $\Kcal_{\mathrm{glob}}=\ker\beta$, fix a puncture
with completed local ring $R$, fraction field $K$, canonical lattice
$\mathfrak p_R$, and completed line $\widehat M$.  Then the completed stalk
is exactly
\[
 \left(N_{\mathfrak{sl}_r(K)}(K\widehat M)\cap\mathfrak p_R\right)/\widehat M.
\]
\end{proposition}

\begin{proof}
The image of $\beta$ is torsion-free, so the kernel is saturated and locally
free.  Completion is flat and commutes with the kernel.  Choose a generator
$m$ of the free rank-one $R$-module $\widehat M$.  For
$x\in\mathfrak p_R$,
\[
 \bar x\in\ker\beta_R
 \Longleftrightarrow [m,x]\in Rm.
\]
The canonical chain-endomorphism lattice is closed under bracket.  Since the
global line is saturated, completion remains saturated and
\[
 (K\widehat M)\cap\mathfrak p_R=\widehat M,
\]
so the preceding condition is equivalent to $[m,x]\in Km$, i.e. to
$x\in N_{\mathfrak{sl}_r(K)}(Km)\cap\mathfrak p_R$.
\end{proof}

In the final high-HN branch, $r=3$ and the universal-extension argument gives
$\Ocal_C^4\subset\Kcal_{\mathrm{glob}}$, generically an equality.  Taking
determinants and applying
\cref{cor:global-integral-det,prop:kernel-compatibility}
gives
\[
 \Ocal_C\longrightarrow\det\Kcal_{\mathrm{glob}}
 \longrightarrow M.
\]
This is impossible because $\deg M=-1$.  Thus the determinant argument
independently recovers the elimination of the final high-HN branch for
arbitrary contact order.

\section{Propagation to all rank-three representations}\label{sec:propagation}

For an irreducible unitary rank-three system, the preceding arguments supply
\[
 H^0(B,R^1\pi_*\ad\Vcal)=0
\]
on every finite etale moduli cover $B$ in every genus $g\ge5$.  The smaller
actual adjoint constituents lie in the published strict range or in the
endpoint lemma.  The Leray argument of
\cite[Proposition 8.2.1]{LLCan} gives strong cohomological rigidity.

The formal constancy step needs vanishing on the dominant etale versal base
used for a linear lift.  The following proposition separates that issue from
the finite-cover geometry.

\begin{proposition}[Projective finite-index descent]
\label{prop:finite-index-descent}
Let $B\to\Mcal_{g,n}$ be a connected finite etale scheme cover carrying the
canonical projective globalization of an irreducible MCG-finite
rank-three representation, and let
\[
 \Ucal_B=\ad(\mathbb P\rho_B)
\]
be its honest adjoint local system.  Suppose that for every connected finite
etale $B'\to B$,
\[
 H^0(B',R^1\pi_*\Ucal_{B'})=0.
 \tag{\ref{prop:finite-index-descent}.1}
\]
Let $M\to\Mcal_{g,n}$ be connected, dominant and etale, and let $\gamma_0$
be a complex local system on the punctured total curve whose fibral
restriction is the same irreducible representation.  Then
\[
 H^0(M,R^1\pi_*\ad\gamma_0)=0.
\]
\end{proposition}

\begin{proof}
Let $X_0$ be the connected component of
$M\times_{\Mcal_{g,n}}B$ containing the chosen matching lift.  It is
surjective finite etale over $M$ and dominant etale over $B$; a nonzero flat
section remains nonzero after pullback.  After one fibral conjugation, the linear local system
$\gamma_0|_{X_0}$ restricts exactly to the chosen irreducible fibre
representation $\rho$.  For an element $x$ of the total fundamental group,
$\gamma_0(x)$ is therefore an exact linear intertwiner from $\rho$ to
$\rho^x$.  By the construction in \cite[Lemma 2.2.2]{LLCan}, the canonical
projective globalization assigns to $x$ the projective class of another
exact linear intertwiner $M_x$ for the same pair.  Thus
$M_x^{-1}\gamma_0(x)$ commutes with $\rho$ and is scalar by Schur's lemma.
Consequently the projective classes agree and
\[
 \ad\gamma_0\simeq f^*\Ucal_B.
\]
The use of exact linear intertwiners is essential: arbitrary projective
extensions may differ by a finite self-twist, which can alternatively be
killed after a further finite etale cover.
The image
\[
 \Gamma=\operatorname{im}(\pi_1(X_0)\to\pi_1(B))
\]
has finite index by \cite[Lemma 2.1.4]{LLCan}.  A global section is therefore
a vector fixed by $\Gamma$.  It becomes a nonzero global section on the
connected finite topological cover corresponding to $\Gamma$, which is a
finite etale algebraic cover $B_\Gamma\to B$ by Riemann existence,
contradicting (\ref{prop:finite-index-descent}.1).
\end{proof}

\begin{lemma}[Direct Artin devissage]\label{lem:artin-devissage}
Let $A$ be a local Artin $\C$-algebra and let $\Wcal_A$ be a local system of
free $A$-modules on a punctured family.  Put
$\Wcal_0=\Wcal_A\otimes_A\C$.  If
\[
 R^0\pi_*\Wcal_0=0,
 \qquad
 H^0(M,R^1\pi_*\Wcal_0)=0,
\]
then
\[
 R^0\pi_*\Wcal_A=0,
 \qquad
 H^0(M,R^1\pi_*\Wcal_A)=0.
\]
\end{lemma}

\begin{proof}
Induct on the length of $A$.  For a small quotient
\[
 0\to I\to A\to A'\to0,
 \qquad \mathfrak m_AI=0,
\]
freeness gives
\[
 0\to I\otimes_\C\Wcal_0\to\Wcal_A\to\Wcal_{A'}\to0.
\]
The action on the kernel is residual.  The long exact sequence of higher
direct images and induction first give $R^0\pi_*\Wcal_A=0$, and then an exact
segment
\[
 0\to I\otimes R^1\pi_*\Wcal_0
 \to R^1\pi_*\Wcal_A
 \to R^1\pi_*\Wcal_{A'}.
\]
Taking global sections proves the claim.
\end{proof}

The preceding argument applies only when the residual unitary fibre is
irreducible.  This qualification is essential: \cite[Lemma 8.5.1]{LLCan}
assumes only unitarity, and the deformation to a PVHS need not preserve
irreducibility.  We now treat every residual rank-three type.

\begin{lemma}[Exact trivial-coefficient gap]
\label{lem:constant-gap}
Let $g\ge2$.  Every nonzero sub-local system of
$R^1\pi_*^\circ\C$ has rank at least $2g$.
\end{lemma}

\begin{proof}
Apply \cite[Lemma 6.1.1]{LLCan} to an irreducible sub-local system $L$.
For the trivial coefficient, $E=\Ocal_C$ and every nonzero
$v\in H^0(\omega_C(D))$ acts injectively on $H^0(\omega_C)$, so the
multiplication rank is $g$.  Lemma 6.1.1 bounds it above by $\rk L/2$.
Therefore $\rk L\ge2g$.
\end{proof}

\begin{lemma}[Reducible rank-three adjoint multiplicities]
\label{lem:reducible-multiplicity}
Let $R$ be semisimple, reducible, non-scalar, and of rank three.  In an
isotypic decomposition
\[
 \ad R\simeq\bigoplus_i U_i^{\oplus w_i},
 \qquad u_i=\rk U_i,
\]
the possible multiplicities satisfy
\[
\begin{array}{c|c}
\text{type of }R&\text{possible }(u_i,w_i)\\ \hline
1+1+1,\text{ not scalar}&(1,w\le4),\\
2+1&(1,w\le3),(2,w\le3),(3,1).
\end{array}
\]
In particular, for $g\ge5$ one has $w_i<2g-2u_i$.
\end{lemma}

\begin{proof}
For three characters, $\ad R$ is the two-dimensional diagonal part together
with the six ratios $\chi_i\chi_j^{-1}$.  If two characters agree, the
trivial coefficient has multiplicity four and each cross character has
multiplicity two; if all three are distinct, every multiplicity is at most
three.

For $R=T\oplus\chi$, with $T$ irreducible of rank two,
\[
 \ad R\simeq\mathbf1\oplus\ad T
 \oplus(T\otimes\chi^{-1})\oplus(T^\vee\otimes\chi).
\]
A rank-one constituent can occur at most three times inside $\ad T$; a
rank-two constituent can occur in the two cross terms and at most once in
$\ad T$; and a rank-three constituent occurs at most once.  At genus five
the corresponding fixed-part lower bounds are $8,6,4$, respectively, and
the inequalities improve with genus.
\end{proof}

\begin{proposition}[Rank-three formal propagation]
\label{prop:formal-propagation}
Assume $g\ge5$, and assume the finite-cover adjoint fixed-part vanishing
proved above for every irreducible unitary MCG-finite rank-three fibre.  Let
$\gamma$ be any free Artin deformation on a dominant etale versal base whose
fibral restriction is the constant deformation of a unitary rank-three
system $R$.  Then
\[
 H^0(M,R^1\pi_*^\circ\ad\gamma)=0.
\]
\end{proposition}

\begin{proof}
If $R$ is irreducible, Schur gives $R^0\pi_*^\circ\ad R=0$,
Proposition~\ref{prop:finite-index-descent} gives residual fixed-part
vanishing, and Lemma~\ref{lem:artin-devissage} propagates it through every
small extension.

Suppose $R$ is reducible and non-scalar.  After the dominant etale base
change of \cite[Lemma 2.4.2]{LLCan}, write
\[
 \ad\gamma\simeq\bigoplus_i U_i\otimes\pi^*W_i,
\]
where $W_i$ is free of rank $w_i$.  The maximal-ideal filtration on $W_i$
has graded pieces that are direct sums of copies of
$W_i^0=W_i\otimes_A\C$.  A first nonzero invariant therefore produces a
nonzero complex map
\[
 (W_i^0)^\vee\longrightarrow R^1\pi_*^\circ U_i
\]
whose image has rank at most $w_i$.  Lemma~\ref{lem:reducible-multiplicity}
gives $w_i<2g-2u_i$, contradicting \cite[Theorem 1.7.1]{LLCan}.

Finally suppose $R=\chi^{\oplus3}$.  Then $\ad R$ is eight copies of the
trivial fibral coefficient.  In the decomposition of
\cite[Lemma 2.4.2]{LLCan}, any rank-one unitary factor that is trivial on
the fibres is pulled back from the base and is absorbed into the coefficient
module.  Write the result as $\ad\gamma\simeq\pi^*W$, where $W$ is free of
rank eight, and put $W^0=W\otimes_A\C$.  A first nonzero Artin graded piece
then gives a nonzero map
\[
 (W^0)^\vee\longrightarrow R^1\pi_*^\circ\C
\]
whose image is a sub-local system of rank at most eight.  This contradicts
Lemma~\ref{lem:constant-gap}, since $8<2g$ for $g\ge5$.
\end{proof}

Proposition~\ref{prop:formal-propagation} supplies exactly the hypothesis of
\cite[Proposition 3.2.1]{LLCan} for every Artin truncation.  It does not
assume that a finite-cover geometric contradiction persists on an arbitrary
dominant etale base.

\paragraph{Published propagation map.}
For clarity, the remaining uses of the strict Landesman--Litt rank bound are
replaced as follows.  In the proof of
\cite[Proposition 8.2.1]{LLCan}, the hypothesis $r<\sqrt{g+1}$ is used to get
$\rk\ad(V)=r^2-1<g$ and invoke \cite[Theorem 6.2.1]{LLCan}; the fixed-part
vanishing and finite-index descent proved above supply that input here.  The
proofs of \cite[Lemmas 8.3.3 and 8.3.4]{LLCan} then use strong rigidity
together with irreducibility, $g\ge3$, quasi-unipotent boundary monodromy,
and finite determinant, but no further rank estimate.  The isomonodromy input
used in \cite[Proposition 8.4.1]{LLCan} has the larger range
$r<2\sqrt{g+1}$, which holds for $r=3$ and $g\ge5$.

In the proof of \cite[Lemma 8.5.1]{LLCan}, the strict bound is again used only
to obtain $\rk\ad(\gamma)=r^2-1<g$ and apply
\cite[Theorem 6.2.1]{LLCan} at every Artin truncation;
Proposition~\ref{prop:formal-propagation} replaces that application.  The
proof of \cite[Lemma 8.5.2]{LLCan} uses the rank bound only through
Lemma~8.5.1.  In the proof of \cite[Lemma 8.6.1]{LLCan}, the original bound
enters in the estimate
\[
 n_i\dim\sigma_i\le \dim\rho_1\dim\rho_2<\frac{g+1}{4}
\]
before \cite[Theorem 7.2.1]{LLCan} is applied.  Here total rank at most three
instead gives $n_i\dim\sigma_i\le2$, while the same published lower bound is
at least six.  Proposition~\ref{prop:rank3-socle} records that replacement,
and no further numerical bound enters the socle induction of
\cite[Section 8.7]{LLCan}.

\begin{proposition}[Unitary rank-three finite image]
\label{prop:unitary-rank3}
Let $g\ge5$.  Every unitary MCG-finite rank-three representation has finite
image.
\end{proposition}

\begin{proof}
If the representation is reducible, its irreducible summands have rank at
most two and have finite image by the published strict theorem.  Assume it is
irreducible.  If its projective image is finite, then its image is finite:
the scalar kernel has finite image because the determinant has finite image,
and the kernel of $z\mapsto z^3$ is finite.  We may therefore assume the
projective image is infinite.  The specialized adjoint fixed-part vanishing and the Leray
argument of \cite[Proposition 8.2.1]{LLCan} give strong cohomological
rigidity.  The proofs of \cite[Lemmas 8.3.3 and 8.3.4]{LLCan} use the
numerical rank hypothesis only through that rigidity; their remaining inputs
are $g\ge3$, finite determinant, and quasi-unipotent boundary monodromy.
Thus the representation is integral.

Strong rigidity is preserved under every complex embedding of its number
field of definition.  Lemma~\cite[Lemma 4.3.2]{LLCan} makes every conjugate
total-space representation a complex PVHS.  The isomonodromy input used in
\cite[Proposition 8.4.1]{LLCan} applies in the larger range
$3<2\sqrt{g+1}$, automatic here, so every conjugate fibre representation is
unitary.  In the current arXiv version of \cite{LLGeo} this input is
Theorem~1.2.13; \cite{LLCan} cites an earlier numbering, Theorem~1.2.12.  The arithmetic unitary criterion at the end of that proposition
then gives finite image.
\end{proof}

\begin{proposition}[Semisimple rank-three finite image]
\label{prop:semisimple-rank3}
Let $g\ge5$.  Every semisimple MCG-finite rank-three representation has
finite image.
\end{proposition}

\begin{proof}
The reducible case has rank-one and rank-two summands and is published.  For
an irreducible system, use the finite-determinant linear globalization and
the constant-determinant nonabelian-Hodge deformation of
\cite[Theorem 4.3.1]{LLCan}.  Its PVHS endpoint is unitary because
$3<2\sqrt{g+1}$ and is MCG-finite by \cite[Proposition 2.1.3]{LLCan}.
Proposition~\ref{prop:unitary-rank3} makes the endpoint finite.
Proposition~\ref{prop:formal-propagation} is precisely the replacement for
\cite[Theorem 6.2.1]{LLCan} in \cite[Lemma 8.5.1]{LLCan}; hence the formal
restriction is constant.  In the proof of \cite[Lemma 8.5.2]{LLCan}, the
numerical rank hypothesis is used only through Lemma 8.5.1.  The remaining
argument spreads formal constancy to a dense open and then uses closedness of
the semisimple conjugacy orbit.  It therefore shows that the original fibre
is conjugate to the unitary endpoint.
\end{proof}

\begin{proposition}[Rank-three characteristic extensions]
\label{prop:rank3-socle}
Let $g\ge5$.  An MCG-finite extension of two semisimple finite-image
representations of total rank at most three has finite image, and hence
splits, provided the subobject is characteristic.
\end{proposition}

\begin{proof}
Let the factor ranks be $a,b$.  Then $a+b\le3$, so $ab\le2$.  Decompose
\[
 \rho_2^\vee\otimes\rho_1
 =\bigoplus_i\sigma_i^{\oplus n_i},
 \qquad d_i=\dim\sigma_i.
\]
Thus $n_id_i\le2$ and $d_i\le2$.  Pass to a finite Galois cover
$\Sigma'_{g',n'}\to\Sigma_{g,n}$ on which both factors are trivial.  The
one-step unipotent restriction factors through
an equivariant homology map
\[
 H_1(\Sigma'_{g',n'};\C)\longrightarrow
 \rho_2^\vee\otimes\rho_1.
\]
Because the subrepresentation is characteristic, the argument of
\cite[Lemma 8.6.1]{LLCan} makes each projected isotypic quotient stable under
a finite-index mapping-class subgroup.  Its dimension is at most
$n_id_i\le2$.  Equivariant duality gives a stable subrepresentation of the
same dimension.  By \cite[Theorem 7.2.1]{LLCan}, every nonzero such
subrepresentation has dimension at least
\[
 2g-2d_i\ge6.
\]
Hence every projected map vanishes and the restriction to the finite cover
is trivial.  Thus the original image is finite; over $\C$ the extension
then splits.
\end{proof}

The semisimplification of an arbitrary MCG-finite rank-three representation
is MCG-finite and finite by Proposition~\ref{prop:semisimple-rank3}.
Applying Proposition~\ref{prop:rank3-socle} inductively to the characteristic
socle filtration, exactly as in \cite[Section 8.7]{LLCan}, proves
\cref{thm:main}.

\section*{Status of the argument}

The theorem is new and unrefereed.  The endpoint and coparabolic convention,
the genus-five branch reconstruction, the normalizer boundary-cocycle
obstruction, and the formal-propagation interfaces are derived explicitly in
the manuscript and cited companion notes.  The spectral-projector closure
and periodic-chain determinant theorem are retained as separate checks,
not as load-bearing inputs.  The included scripts verify finite arithmetic
and fibrewise linear algebra; they do not verify the Hodge, parabolic, or
deformation interfaces imported from Landesman--Litt.

\paragraph{Assistance disclosure.}
OpenAI Codex (GPT-5.6) assisted with preparation and internal checking of this
manuscript.  The author reviewed the manuscript and is responsible for all
statements, proofs, and errors.

\begin{thebibliography}{99}

\bibitem{BPGN}
L.~Brambila-Paz, I.~Grzegorczyk, and P.~E. Newstead,
\emph{Geography of Brill--Noether loci for small slopes},
J. Algebraic Geom. \textbf{6} (1997), 645--669.

\bibitem{ChenSalter}
L.~Chen and N.~Salter,
\emph{The Birman exact sequence does not virtually split},
Math. Res. Lett. \textbf{28} (2021), 383--413.

\bibitem{EarleKra}
C.~J. Earle and I.~Kra,
\emph{On sections of some holomorphic families of closed Riemann surfaces},
Acta Math. \textbf{137} (1976), 49--79.

\bibitem{LLCan}
A.~Landesman and D.~Litt,
\emph{Canonical representations of surface groups},
Ann. of Math. (2) \textbf{199} (2024), 823--897.

\bibitem{LLGeo}
A.~Landesman and D.~Litt,
\emph{Geometric local systems on very general curves and isomonodromy},
J. Amer. Math. Soc. \textbf{37} (2024), 683--729.

\end{thebibliography}

\end{document}
